\documentclass[11pt]{article}

\usepackage[T1]{fontenc}
\usepackage[utf8]{inputenc}
\usepackage{lmodern}
\usepackage[margin=1.1in]{geometry}
\usepackage{amsmath,amssymb,amsthm,mathtools}
\usepackage{enumitem}
\usepackage{microtype}
\usepackage{xcolor}
\usepackage{booktabs}
\usepackage{array}
\usepackage{titlesec}
\usepackage{hyperref}

\titleformat{\section}
  {\normalfont\LARGE\bfseries}{\thesection}{1em}{}
\titleformat{\subsection}
  {\normalfont\Large\bfseries}{\thesubsection}{1em}{}

\hypersetup{
  pdftitle={Collision Ideals and Off-Diagonal Sheets},
  pdfauthor={Chloe van der Vlugt},
  colorlinks=true,
  linkcolor=blue!55!black,
  citecolor=green!40!black,
  urlcolor=blue!60!black
}

\newtheorem{theorem}{Theorem}[section]
\newtheorem{proposition}[theorem]{Proposition}
\newtheorem{lemma}[theorem]{Lemma}
\newtheorem{corollary}[theorem]{Corollary}
\theoremstyle{remark}
\newtheorem*{remark}{Remark}
\theoremstyle{definition}
\newtheorem{definition}[theorem]{Definition}

\DeclareMathOperator{\Spec}{Spec}
\DeclareMathOperator{\Frac}{Frac}
\DeclareMathOperator{\Gal}{Gal}
\DeclareMathOperator{\Hom}{Hom}
\DeclareMathOperator{\Aut}{Aut}
\DeclareMathOperator{\Ram}{Ram}
\DeclareMathOperator{\Norm}{Norm}
\DeclareMathOperator{\Ann}{Ann}
\DeclareMathOperator{\core}{core}

\newcommand{\A}{\mathbb A}
\newcommand{\C}{\mathbb C}
\newcommand{\Obs}{\operatorname{Obs}}
\newcommand{\id}{\mathrm{id}}
\newcommand{\qedbox}{\hfill$\square$}
\newcommand{\leanpath}[1]{\texttt{\nolinkurl{#1}}}

\newenvironment{leanfiles}
  {\par\medskip\begingroup\footnotesize\raggedright\noindent
   \textit{Lean correspondence.}\ }
  {\par\endgroup\medskip}

\title{Collision Ideals and Off-Diagonal Sheets}
\author{Chloe van der Vlugt}
\date{\today}

\begin{document}

\maketitle

\begin{abstract}
We study the scheme-theoretic collision relation of a polynomial self-map,
namely its fiber product over the target, and describe the geometry away
from the diagonal through boundary-supported inertia and Galois-closure
sheets.
For a polynomial map
\(F=(F_1,\dots,F_n)\colon\A^n_{\C}\to\A^n_{\C}\), define in
\(S=\C[x_1,\dots,x_n,y_1,\dots,y_n]\) the collision and diagonal ideals
\[
  I_R=\bigl(F_i(x)-F_i(y)\bigr)_{i=1}^n,
  \qquad
  I_\Delta=(x_i-y_i)_{i=1}^n.
\]
The image of \(I_\Delta\) in the collision algebra \(C_F:=S/I_R\) is the
obstruction ideal
\(\Obs(F):=I_\Delta/I_R\); it records the failure of the collision scheme
\(\Spec C_F\) to coincide with the diagonal.  Its vanishing characterizes
polynomial automorphisms:
\[
  \Obs(F)=0
  \quad\Longleftrightarrow\quad
  F\text{ is a polynomial automorphism}.
\]
In dimension two we give two module-theoretic formulations of the same
fixed-map conclusion.  Planar boundary coherence asks that the boundary
local cohomology
\(H^1_{\partial X}(\overline X,\mathcal O_{\overline X})\) be a finite
\(\overline A\)-module; it forces \(\partial X=\varnothing\), hence makes
the Keller map finite \'etale and therefore an automorphism.  Planar
ramification rigidity is the divisorial criterion that
\(\Omega_{A_N/B}\) have no height-one support, equivalently that it have
finite length.  On the common Galois normalization \(Z=\Spec A_N\), we
construct the pulled-back
affine opens \(U_g\), their canonical boundary ideals
\(\mathfrak j_g=I_{A_N}(Z\setminus U_g)\), and the fixed-locus ideals
\(\mathfrak a_C\).  The height-one dictionary identifies moving-sheet
coverage and finite-group boundary separation as exact criteria for
ramification rigidity, after which purity and finite-\'etale rigidity give
\(N=L=K\) and collision vanishing.  Although \(\overline A\) and \(A_N\)
are already finite over \(B\), neither of these additional finiteness
criteria follows formally.  Universal validity of either criterion is
equivalent to the two-dimensional Jacobian conjecture and is not proved
here.  The boundary divisor-class argument does, however, prove the
normal/Galois case and therefore excludes generic degree two.
For the separable nonnormal cubic class in dimension three, let
\[
  K=\C(F_1,F_2,F_3),
  \qquad
  L=\C(x_1,x_2,x_3),
\]
and let \(N\) be the normal closure of \(L/K\).  Then
\[
  L\otimes_K L\simeq L\times N,
  \qquad
  \Gal(N/K)\simeq S_3.
\]
The generic off-diagonal factor \(N\) is the ordered-root Galois-closure sheet.
Consequently, \(\Obs(F)\neq0\).  The algebraic construction, the
pulled-back boundary dictionary, and the implications from supplied
coverage or separation hypotheses are formalized in Lean~4.  The
boundary-coherence and finite-length support arguments remain explicit
Lean interfaces whose manuscript proofs have not yet been formalized, as
do the standard external geometric statements.
\end{abstract}

\begin{center}
  \small Licensed under
  \href{https://creativecommons.org/licenses/by/4.0/}
    {CC BY 4.0}.
\end{center}

\tableofcontents

\section{Introduction}
\label{sec:introduction}

The Jacobian conjecture asks whether a polynomial self-map of affine space
with constant nonzero Jacobian determinant must be a polynomial
automorphism.  Dimension two remains open.  Rather than treating
injectivity as an external property of the map, this paper studies the
scheme-theoretic fiber product of the map with itself and asks when its
off-diagonal part disappears.

There is one divisorial endgame, expressed in several exact languages.  The
paper does not prove the Jacobian conjecture.  It organizes the remaining
problem around the following two module-theoretic formulations:
\[
  \begin{aligned}
  \operatorname{PBC}(F)
    &:\Longleftrightarrow
      H^1_{\partial X}(\overline X,\mathcal O_{\overline X})
      \text{ is a finite }\overline A\text{-module},\\
  \operatorname{PRR}(F)
    &:\Longleftrightarrow
      \Omega_{A_N/B}\text{ has no height-one support}\notag\\[-2pt]
    &\hspace{4.7em}\Longleftrightarrow
      \Omega_{A_N/B}\text{ has finite length over }A_N.
  \end{aligned}
\]
Planar boundary coherence \(\operatorname{PBC}(F)\) is the corresponding
global properness formulation:
\[
  \operatorname{PBC}(F)
  \Longrightarrow \partial X=\varnothing
  \Longrightarrow f_F\text{ finite \'etale}
  \Longrightarrow F\text{ an automorphism}.
\]
Planar ramification rigidity \(\operatorname{PRR}(F)\) is the
differential-module formulation of the codimension-one premise:
\[
  \operatorname{PRR}(F)
  \Longrightarrow Z/Y\text{ finite \'etale}
  \Longrightarrow N=K
  \Longrightarrow L=K
  \Longrightarrow q_F=0
  \Longrightarrow I_R=I_\Delta.
\]

The divisorial endgame is controlled by an exact codimension-one obstruction on
the Galois normalization.  A ramified divisor \(E\) is eliminated as soon
as one conjugate sheet moved by \(I_E\) still contains its generic point.
In coordinates this is the conjugate-integrality condition
\[
  I_E\neq1
  \Longrightarrow
  \exists g\in G,\quad
  I_E\nsubseteq gHg^{-1},\qquad
  w_E(g(x_1)),w_E(g(x_2))\geq0.
\]
We call this planar moving-sheet coverage.  Its ideal-theoretic form says
that the fixed locus of a prime-order subgroup and the boundaries of all
the sheets it moves have no common height-one component.  Establishing
this universally---equivalently, establishing planar ramification rigidity
universally---is equivalent to \(JC(2)\), so the reformulation locates the
missing theorem without assuming it.

Only after all height-one points have been shown unramified do purity and
finite-\'etale rigidity enter: they make the normal closure trivial and
collapse the collision relation to the diagonal.  This order is essential
to the argument.  The paper also develops boundary divisor-class and
principal-parts criteria that may imply coverage under additional
hypotheses; these are prospective mechanisms, not additional endgames or
proofs of universal coverage.
The divisor-class mechanism does prove unconditionally that a normal---hence
Galois---marked extension \(L/K\) is trivial.  In particular, a hypothetical
planar counterexample must have \(L/K\) nonnormal and generic degree at least
three.

The normalization rings \(\overline A\) and \(A_N\) are finite
\(B\)-modules from the outset.  This makes \(\overline X\) and \(Z\)
affine varieties: they are integral affine schemes of finite type over
\(\C\).  It does not make \(A\) finite over \(B\), make boundary local
cohomology finite, or make \(\Omega_{A_N/B}\) finite length.  Likewise,
the vanishing of ordinary higher quasi-coherent cohomology on an affine
scheme does not imply the vanishing or finiteness of local cohomology with
support in its deleted boundary.

In the separable nonnormal cubic class in dimension three, the opposite
phenomenon is visible generically: the off-diagonal factor is the
ordered-root \(S_3\)-normal-closure sheet.  This gives a structural
comparison with the planar criteria while keeping its assumptions explicit.

Section~\ref{sec:general-construction} develops collision ideals and the
dimension-independent normalization diagram.
Section~\ref{sec:planar-vanishing} constructs the planar pulled-back
boundaries, proves the valuation dictionary, and states the equivalent
coverage and separation criteria before applying the endgame.
Section~\ref{sec:cubic-s3-obstruction} treats the cubic \(S_3\) class.

\section{General construction}
\label{sec:general-construction}

The organizing question is how the failure of a polynomial self-map to be
one-to-one is encoded by the part of its fiber square lying off the diagonal.
The collision ideal first isolates that part scheme-theoretically; the
normalization diagram then describes it through conjugate sheets and inertia.
These constructions are dimension-independent.  Section
\ref{sec:planar-vanishing} specializes them to planar Keller maps and locates
the obstruction at the deleted normalization boundary, while Section
\ref{sec:cubic-s3-obstruction} identifies the generic off-diagonal sheet in
the separable nonnormal cubic class with an \(S_3\)-normal closure.  The two
specializations respectively describe when the collision scheme collapses to
the diagonal and what its surviving generic sheet looks like when it does not.

\subsection{The two ideals and their obstruction}
\label{sec:collision-ideals}

Throughout the paper, the base field is \(\C\).  Let
\[
  F=(F_1,\dots,F_n)\colon
  X=\A^n_{\C}\longrightarrow\A^n_{\C}
\]
be a polynomial map, and let
\[
  S=\C[x_1,\dots,x_n,y_1,\dots,y_n].
\]
The map \(F\) is Keller if
\[
  \det JF\in\C^\times.
\]
Define
\[
  I_R(F)=\bigl(F_i(x)-F_i(y)\bigr)_{i=1}^{n},
  \qquad
  I_\Delta=(x_i-y_i)_{i=1}^{n},
  \qquad
  \Obs(F):=I_\Delta/I_R(F).
\]
The quotient \(\Obs(F)\) is naturally an ideal of the collision ring
\(C_F:=S/I_R(F)\); we call it the obstruction ideal.
For a fixed map \(F\), we abbreviate \(I_R(F)\) to \(I_R\).

\begin{proposition}[Canonical containment]
\label{prop:collision-contained-in-diagonal}
For every polynomial map \(F\),
\[
  I_R(F)\subseteq I_\Delta.
\]
\end{proposition}

\begin{proof}
For \(H\in\C[x_1,\dots,x_n]\), the difference \(H(x)-H(y)\) belongs to
\((x_1-y_1,\dots,x_n-y_n)\).  This follows by induction on \(H\), or by
telescoping each monomial one coordinate at a time.  Apply the statement
to each coordinate \(F_i\).
\end{proof}

Let \(A:=\C[x_1,\dots,x_n]\).  Diagonal evaluation induces a surjective ring
map
\[
  \bar\mu_F\colon C_F\longrightarrow A,
  \qquad
  [s]_{I_R}\longmapsto s(x,x).
\]

\begin{proposition}[The diagonal kernel]
\label{prop:obstruction-kernel}
As ideals of \(C_F\),
\[
  \ker(\bar\mu_F)
  =
  \Obs(F)
  =
  I_\Delta/I_R.
\]
Consequently,
\[
  \ker(\bar\mu_F)=0
  \quad\Longleftrightarrow\quad
  \Obs(F)=0
  \quad\Longleftrightarrow\quad
  I_R=I_\Delta.
\]
\end{proposition}

\begin{proof}
The kernel of diagonal evaluation \(S\to A\) is \(I_\Delta\).  Passing
through \(S/I_R\) therefore identifies the displayed kernel with
\(I_\Delta/I_R\).  Since \(I_R\subseteq I_\Delta\),
\[
  I_\Delta/I_R=0
  \quad\Longleftrightarrow\quad
  I_R=I_\Delta.
\]
\end{proof}

\subsection{The fiber-product interpretation}
\label{sec:fiber-product}

Let
\[
  A=\C[x_1,\dots,x_n],\qquad
  B=\C[F_1,\dots,F_n]\subseteq A.
\]
Let \(X=\Spec A\), let \(Y=\Spec B\), and let
\[
  f_F\colon X\longrightarrow Y
\]
be the morphism induced by \(B\hookrightarrow A\).
The homomorphisms
\(\iota_x,\iota_y\colon A\rightrightarrows C_F\), defined by
\(\iota_x(a)=[a(x)]\) and \(\iota_y(a)=[a(y)]\), agree on \(B\) because
\([F_i(x)]=[F_i(y)]\) in \(C_F\).  Their common restriction equips
\(C_F\) with a \(B\)-algebra structure.

\begin{theorem}[Fiber-product interpretation]
\label{thm:collision-tensor-product}
There is a canonical \(B\)-algebra equivalence and its dual
affine-scheme equivalence
\[
  C_F\cong A\otimes_B A,
  \qquad
  R_F:=\Spec C_F\cong X\times_YX.
\]
Under the first equivalence, \(\bar\mu_F\) is tensor multiplication
\[
  A\otimes_BA\longrightarrow A,\qquad a\otimes a'\longmapsto aa'.
\]
Its dual morphism is the diagonal
\[
  \Delta_{X/Y}\colon X\longrightarrow X\times_YX.
\]
\end{theorem}

\begin{proof}
A map \(C_F\to T\) to a commutative \(\C\)-algebra \(T\) is equivalent
to a pair of maps \(f,g\colon A\rightrightarrows T\) satisfying
\[
  f(F_i)=g(F_i)\qquad(1\leq i\leq n).
\]
Since the \(F_i\) generate \(B\), this is precisely the universal
property of \(A\otimes_BA\).  Applying \(\Spec\) to this equivalence and
using
\[
  \Spec(A\otimes_BA)
  \cong
  \Spec A\times_{\Spec B}\Spec A
\]
gives the asserted fiber-product equivalence.
The statements about multiplication and the diagonal follow on pure
tensors.
\end{proof}

\[
  R_F(\C)
  =
  \{(u,v)\in X(\C)\times X(\C):F(u)=F(v)\}.
\]
Inside
\[
  X\times X=\Spec S,
\]
the ideals \(I_R\) and \(I_\Delta\) define, respectively,
\(X\times_YX\) and the image of \(\Delta_{X/Y}\).  The
contravariance of \(I\mapsto V(I)\) therefore gives the equivalences
\[
  I_R\subseteq I_\Delta
  \quad\Longleftrightarrow\quad
  \Delta_{X/Y}(X)\subseteq X\times_YX
\]
and
\[
  I_R=I_\Delta
  \quad\Longleftrightarrow\quad
  X\times_YX=\Delta_{X/Y}(X)
\]
as closed subschemes of \(X\times X\).  Thus the algebraic and
geometric gaps are two descriptions of the same scheme-theoretic
containment.

\subsection{The off-diagonal factor and secant projectors}
\label{sec:secant-projector}

For every polynomial map \(F\), define
\[
  I_{\mathrm{off}}:=I_R:I_\Delta,
  \qquad
  C_F^\circ:=S/I_{\mathrm{off}},
  \qquad
  R_F^\circ:=\Spec C_F^\circ.
\]
This definition does not require the diagonal to be clopen.

\begin{corollary}[Off-diagonal emptiness criterion]
\label{cor:off-diagonal-empty-criterion}
For every polynomial map \(F\), the following are equivalent:
\[
  R_F^\circ=\varnothing
  \quad\Longleftrightarrow\quad
  C_F^\circ=0
  \quad\Longleftrightarrow\quad
  \Obs(F)=0
  \quad\Longleftrightarrow\quad
  I_R=I_\Delta.
\]
\end{corollary}

\begin{proof}
\[
  R_F^\circ=\varnothing
  \quad\Longleftrightarrow\quad
  C_F^\circ=0
  \quad\Longleftrightarrow\quad
  I_R:I_\Delta=S.
\]
Moreover,
\[
  I_R:I_\Delta=S
  \quad\Longleftrightarrow\quad
  I_\Delta\subseteq I_R
  \quad\Longleftrightarrow\quad
  I_R=I_\Delta.
\]
The equivalence with \(\Obs(F)=0\) is
Proposition~\ref{prop:obstruction-kernel}.
\end{proof}

The idempotent mechanism behind a clopen diagonal is purely
commutative-algebraic.  Let \(C\) be a commutative ring, \(J\subseteq C\)
an ideal, \(\delta\in C\), and \(c\in C^\times\).  Suppose
\[
  \delta J=0,
  \qquad
  \delta-c\in J.
\]
Define
\[
  q:=1-c^{-1}\delta.
\]

\begin{theorem}[Abstract off-diagonal projector]
\label{thm:abstract-off-diagonal-projector}
The element \(q\) is idempotent and
\[
  J=Cq,
  \qquad
  \Ann_C(J)=C(1-q).
\]
In particular,
\[
  J=0\quad\Longleftrightarrow\quad q=0.
\]
\end{theorem}

\begin{proof}
The congruence \(\delta\equiv c\pmod J\) gives \(q\in J\).  For every
\(j\in J\),
\[
  qj=j-c^{-1}\delta j=j.
\]
Taking \(j=q\) proves \(q^2=q\); the same identity gives \(J=Cq\).
Finally,
\[
  aq=0
  \quad\Longleftrightarrow\quad
  a=a(1-q),
\]
so \(\Ann_C(J)=\Ann_C(Cq)=C(1-q)\).
\end{proof}

Whenever the theorem is applied with
\[
  C=C_F,
  \qquad
  J=\Obs(F),
\]
the identities \(J=C_Fq\) and \(C_F/J\cong A\) give
\[
  C_F
  \cong
  A
  \times
  C_F/C_F(1-q).
\]
Thus the first factor is the diagonal and the second is its
off-diagonal complement.

\subsection{Generic fibers and marked-root decomposition}
\label{sec:generic-galois-bridge}

Let \(F\) be dominant and generically finite, and let
\[
  A=\C[x_1,\dots,x_n],
  \qquad
  B=\C[F_1,\dots,F_n]\subseteq A.
\]
Let
\[
  K=\Frac(B),
  \qquad
  L=\Frac(A).
\]
Let
\[
  \theta\colon K\otimes_BA\longrightarrow L,
  \qquad
  \lambda\otimes a\longmapsto\lambda a.
\]

\begin{theorem}[Generic collision bridge]
\label{thm:generic-collision-base-change}
The map \(\theta\) is an isomorphism.  Consequently, there is a
\(K\)-algebra isomorphism
\[
  K\otimes_BC_F
  \cong
  L\otimes_KL,
\]
and the following square commutes:
\[
  \begin{array}{ccc}
    K\otimes_BC_F
      &\xrightarrow{\ 1\otimes\bar\mu_F\ }&
    K\otimes_BA\\
    {\scriptstyle\sim}\downarrow
      &&\downarrow{\scriptstyle\theta\ \sim}\\
    L\otimes_KL
      &\xrightarrow{\ \mu_L\ }&
    L.
  \end{array}
\]
\end{theorem}

\begin{proof}
Let \(U=B\setminus\{0\}\).  Since \(K=U^{-1}B\), localization gives
\[
  K\otimes_BA
  \xrightarrow{\sim}
  U^{-1}A,
  \qquad
  \frac{b}{u}\otimes a\longmapsto\frac{ba}{u}.
\]
Since \(A=\C[x_1,\dots,x_n]\) is a domain and \(B\subseteq A\),
\[
  B\setminus\{0\}\subseteq A\setminus\{0\}.
\]
Hence
\[
  U^{-1}A\hookrightarrow\Frac(A)=L,
  \qquad
  \frac{a}{u}\longmapsto\frac{a}{u}.
\]
Composing this injection with the displayed isomorphism gives
\(\theta\), so \(\theta\) is injective.

Since \(F\) is generically finite, \(L/K\) is a finite extension.  The
image of \(U^{-1}A\) in \(L\) is
\[
  K[x_1,\dots,x_n].
\]
Since each \(x_i\) is algebraic over \(K\), this \(K\)-domain is
finite-dimensional over \(K\), hence is a field.  Since it contains
\(A\), it contains \(\Frac(A)=L\); the reverse inclusion is already
given by its embedding into \(L\).  Therefore
\[
  U^{-1}A=L,
\]
and \(\theta\) is an isomorphism.

By Theorem~\ref{thm:collision-tensor-product}, identify
\(C_F\) with \(A\otimes_BA\).  There is a \(K\)-algebra isomorphism
\[
  \beta\colon
  K\otimes_B(A\otimes_BA)
  \xrightarrow{\sim}
  (K\otimes_BA)\otimes_K(K\otimes_BA)
\]
defined on pure tensors by
\[
  \beta\bigl(\lambda\otimes(a\otimes a')\bigr)
  =
  (\lambda\otimes a)\otimes(1\otimes a').
\]
Its inverse is
\[
  (\lambda\otimes a)\otimes(\lambda'\otimes a')
  \longmapsto
  \lambda\lambda'\otimes(a\otimes a').
\]
The balancing relations over \(B\) and \(K\) show that these formulas
are well-defined, and the two displayed maps are mutually inverse.
Composing \(\beta\) with \(\theta\otimes_K\theta\) gives
\[
  \Phi_F\colon K\otimes_BC_F\xrightarrow{\sim}L\otimes_KL.
\]

It remains to identify the diagonal map.  Under
\(C_F\cong A\otimes_BA\), the homomorphism \(\bar\mu_F\) is
multiplication.  Therefore, on a pure tensor,
\[
  \begin{aligned}
  (\mu_L\circ\Phi_F)
    \bigl(\lambda\otimes(a\otimes a')\bigr)
    &=(\lambda a)a'
     =\lambda aa',\\
  \bigl(\theta\circ(1\otimes\bar\mu_F)\bigr)
    \bigl(\lambda\otimes(a\otimes a')\bigr)
    &=\theta(\lambda\otimes aa')
     =\lambda aa'.
  \end{aligned}
\]
Pure tensors generate \(K\otimes_BC_F\), so the two homomorphisms are
equal.  This proves the commutative square.
\end{proof}

Suppose now that \(L/K\) is separable and has a power basis
\(L=K(\alpha)\).  Let \(m_\alpha(T)\in K[T]\) be the minimal polynomial
of \(\alpha\), and let \(h_\alpha(T)\in L[T]\) be determined by
\[
  m_\alpha(T)=(T-\alpha)h_\alpha(T)
  \qquad\text{in }L[T].
\]

\begin{theorem}[Marked-root tensor decomposition]
\label{thm:marked-root-tensor-decomposition}
Give \(L\otimes_KL\) its \(L\)-algebra structure through the left tensor
factor.  There is an \(L\)-algebra equivalence
\[
  L\otimes_KL
  \cong
  L\times L[T]/(h_\alpha),
\]
under which tensor multiplication \(L\otimes_KL\to L\) is projection
onto the first factor.
\end{theorem}

\begin{proof}
The power-basis presentation
\[
  L\cong K[T]/(m_\alpha)
\]
and base change from \(K\) to \(L\) give an \(L\)-algebra isomorphism
\[
  L\otimes_KL
  \xrightarrow{\sim}
  L[T]/(m_\alpha),
  \qquad
  1\otimes\alpha\longmapsto T.
\]
Since \(m_\alpha\) is separable,
\(m_\alpha'(\alpha)\neq0\).  Differentiating
\[
  m_\alpha(T)=(T-\alpha)h_\alpha(T)
\]
and evaluating at \(T=\alpha\) gives
\[
  h_\alpha(\alpha)=m_\alpha'(\alpha)\neq0.
\]
Thus \(T-\alpha\) does not divide \(h_\alpha\), so these two polynomials
are coprime in \(L[T]\).  The Chinese remainder theorem now gives
\[
  L[T]/(m_\alpha)
  \xrightarrow{\sim}
  L[T]/(T-\alpha)\times L[T]/(h_\alpha)
  \xrightarrow{\sim}
  L\times L[T]/(h_\alpha).
\]
Under the first isomorphism above, a pure tensor
\(\ell\otimes p(\alpha)\), with \(p\in K[T]\), corresponds to the class
of \(\ell p(T)\).  Tensor multiplication sends it to
\(\ell p(\alpha)\), which is evaluation at \(T=\alpha\).  Under the
Chinese remainder isomorphism, this evaluation map is projection onto
the factor \(L[T]/(T-\alpha)\cong L\).  Hence tensor multiplication is
the first projection.
\end{proof}

Together with Theorem~\ref{thm:generic-collision-base-change}, this gives
\[
  K\otimes_BC_F
  \cong
  L\times L[T]/(h_\alpha),
\]
where the first projection is the generic diagonal.

\subsection{Normalized covers and inertia}
\label{sec:normalized-covers-inertia}

Let \(L/K\) be finite and separable, and let \(N/K\) be its normal
closure, with a fixed embedding \(L\hookrightarrow N\).  Let
\[
  G=\Gal(N/K),
  \qquad
  H=\Gal(N/L).
\]

\begin{proposition}[Conjugate-embedding action]
\label{prop:conjugate-embedding-action}
Restriction induces a bijection
\[
  G/H\xrightarrow{\sim}\Hom_K(L,N),
  \qquad
  gH\longmapsto g|_L.
\]
The induced action of \(G\) on \(G/H\) is faithful; equivalently,
\[
  \core_G(H)=\bigcap_{g\in G}gHg^{-1}=1.
\]
\end{proposition}

\begin{proof}
For \(g_1,g_2\in G\),
\[
  g_1|_L=g_2|_L
  \quad\Longleftrightarrow\quad
  g_2^{-1}g_1\in H
  \quad\Longleftrightarrow\quad
  g_1H=g_2H.
\]
Every \(K\)-embedding \(L\hookrightarrow N\) extends to an element of
\(G\), so restriction gives the stated bijection.  Its action kernel is
\[
  \ker\!\left(G\curvearrowright G/H\right)
  =
  \bigcap_{g\in G}gHg^{-1}
  =
  \core_G(H).
\]
If \(\sigma\in\core_G(H)\), then
\[
  \sigma|_{g(L)}=\id_{g(L)}
  \quad(g\in G).
\]
Since \(N\) is the normal closure,
\[
  N=K\bigl(g(L):g\in G\bigr),
\]
and therefore \(\sigma=\id_N\).
\end{proof}

Let
\[
  Y=\Spec B,
  \qquad
  \overline X=\Norm_L(Y),
  \qquad
  Z=\Norm_N(Y).
\]
The marked embedding \(L\hookrightarrow N\) induces a canonical
commutative triangle
\[
  Z\xrightarrow{\pi}\overline X\xrightarrow{\bar\nu}Y,
  \qquad
  \nu=\bar\nu\circ\pi\colon Z\longrightarrow Y.
\]
Suppose that \(\bar\nu\) and \(\nu\) are finite, and fix an open
immersion over \(Y\)
\[
  j\colon X=\Spec A\hookrightarrow\overline X
\]
with \(\bar\nu\circ j=f_F\).  Denote this diagram and its deleted boundary
by
\[
  \mathcal D=
  \left(
    Z\xrightarrow{\pi}\overline X\xrightarrow{\bar\nu}Y,\,
    j\colon X\hookrightarrow\overline X
  \right),
  \qquad
  \partial X=\overline X\setminus j(X).
\]

Fix \(E\in Z^{(1)}\).  Let \(w_E\) be its divisorial
valuation, let \(b=\nu(E)\), and let \(v_b=w_E|_K\).
Let \(D_E\leq G\) and \(I_E\trianglelefteq D_E\) be its decomposition and
inertia groups.  Let
\[
  \mathcal Q_E:=D_E\backslash G/H.
\]
For \(\omega=D_EgH\in\mathcal Q_E\), let
\[
  u_\omega=(g^{-1}w_E)|_L,
  \qquad
  \operatorname{cent}(E,\omega)
  =\operatorname{cent}_{\overline X}(u_\omega),
\]
and let
\[
  i_E(\omega)=[I_E:I_E\cap gHg^{-1}].
\]
\begin{proposition}[Conjugate divisors and double cosets]
\label{prop:conjugate-divisor-double-cosets}
The assignments above are well-defined, and restriction of valuations
induces a bijection
\[
  \mathcal Q_E
  \xrightarrow{\sim}
  \{\text{prime divisors of }\overline X\text{ above }b\},
  \qquad
  \omega\longmapsto\operatorname{cent}(E,\omega).
\]
\end{proposition}

\begin{proof}
The group \(G\) acts transitively on the extensions of \(v_b\) to \(N\),
and the stabilizer of \(w_E\) is \(D_E\).  Restricting such a valuation
to \(L=N^H\) identifies precisely those extensions whose indices differ
on the right by an element of \(H\).  The resulting orbits are therefore
indexed by \(D_E\backslash G/H\), giving the stated bijection.

It remains to check independence of representatives.  If
\(D_Eg'H=D_EgH\), then \(g'=dgh\) for some \(d\in D_E\) and \(h\in H\).
The element \(d\) fixes \(w_E\), while \(h\) fixes \(L\) pointwise, so
\[
  (g'^{-1}w_E)|_L=(g^{-1}w_E)|_L.
\]
Moreover, \(I_E\trianglelefteq D_E\), and conjugation by \(d\) identifies
\[
  I_E\cap g'Hg'^{-1}
  \quad\text{with}\quad
  I_E\cap gHg^{-1}.
\]
Thus \(u_\omega\), \(\operatorname{cent}(E,\omega)\), and
\(i_E(\omega)\) are
well-defined.
\end{proof}

We use the notation
\[
  e\bigl(\operatorname{cent}(E,\omega)/b\bigr):=e(u_\omega/v_b).
\]
Since the residue characteristic is zero,
\[
  \Ram^{(1)}(Z/Y)
  :=
  \{E\in Z^{(1)}:e(w_E/v_b)>1\}
  =
  \{E\in Z^{(1)}:I_E\neq1\}.
\]

\begin{proposition}[Detection of inertia on a conjugate sheet]
\label{prop:core-free-sheet-action}
If \(I_E\neq1\), then there is an \(\omega\in\mathcal Q_E\) with
\[
  i_E(\omega)>1.
\]
\end{proposition}

\begin{proof}
Suppose, to the contrary, that \(i_E(\omega)=1\) for every
\(\omega\in\mathcal Q_E\).  For each \(g\in G\),
\[
  1=i_E(D_EgH)=[I_E:I_E\cap gHg^{-1}],
\]
so
\[
  I_E=I_E\cap gHg^{-1}\subseteq gHg^{-1}
  \qquad\text{for every }g\in G.
\]
Consequently,
\[
  I_E
  \subseteq
  \bigcap_{g\in G}gHg^{-1}
  =
  \core_G(H)
  =
  1,
\]
contrary to \(I_E\neq1\).
\end{proof}

\begin{proposition}[Intermediate ramification index]
\label{prop:intermediate-ramification-index}
For \(\omega=D_EgH\in\mathcal Q_E\),
\[
  e\bigl(\operatorname{cent}(E,\omega)/b\bigr)
  =
  i_E(\omega)
  =
  [I_E:I_E\cap gHg^{-1}].
\]
\end{proposition}

\begin{proof}
Let \(w_E\) and \(v_b\) be the divisorial valuations associated with
\(E\) and \(b\), and let
\[
  w_g=g^{-1}w_E,
  \qquad
  u_\omega=w_g|_L.
\]
The valuation \(u_\omega\) has center \(\operatorname{cent}(E,\omega)\) on
\(\overline X\).  Since
the residue fields have characteristic zero,
\[
  e(w_g/v_b)=|I_E|.
\]
The inertia group of \(w_g\) over \(u_\omega\) is
\[
  g^{-1}I_Eg\cap H,
\]
and hence
\[
  e(w_g/u_\omega)
  =
  |g^{-1}I_Eg\cap H|
  =
  |I_E\cap gHg^{-1}|.
\]
Multiplicativity of ramification indices in the field tower
\[
  K\subseteq L\subseteq N
\]
therefore gives
\[
  e(u_\omega/v_b)
  =
  \frac{e(w_g/v_b)}{e(w_g/u_\omega)}
  =
  [I_E:I_E\cap gHg^{-1}],
\]
which is the asserted formula.
\end{proof}

\begin{proposition}[Visibility-to-boundary]
\label{prop:visibility-to-boundary}
Suppose that \(X\to Y\) is \'etale.  For \(E\in Z^{(1)}\) and
\(\omega\in\mathcal Q_E\),
\[
  i_E(\omega)>1
  \quad\Longrightarrow\quad
  \operatorname{cent}(E,\omega)\in\partial X.
\]
\end{proposition}

\begin{proof}
Suppose that \(\operatorname{cent}(E,\omega)\in j(X)\), and let
\(x\in X\) be the corresponding point under the open immersion \(j\).
The local homomorphism
\[
  \mathcal O_{Y,b}
  \longrightarrow
  \mathcal O_{X,x}
\]
is \'etale, so
\(e(\operatorname{cent}(E,\omega)/b)=1\).  Proposition
\ref{prop:intermediate-ramification-index} now gives
\[
  i_E(\omega)
  =
  e\bigl(\operatorname{cent}(E,\omega)/b\bigr)
  =
  1.
\]
Taking the contrapositive gives the result.
\end{proof}

\begin{definition}[Hidden-inertia locus]
\label{def:no-hidden-inertia}
\[
  \mathcal H(\mathcal D)
  :=
  \left\{
  E\in\Ram^{(1)}(Z/Y):
  \begin{array}{l}
    \exists\omega\in\mathcal Q_E,\ i_E(\omega)>1,\\
    \forall\omega\in\mathcal Q_E,\quad
      i_E(\omega)>1
      \Longrightarrow
      \operatorname{cent}(E,\omega)\in\partial X
  \end{array}
  \right\}.
\]
Thus \(E\in\mathcal H(\mathcal D)\) precisely when nontrivial inertia
is visible on a conjugate sheet but every visible center belongs to the
deleted boundary.
\end{definition}

\begin{proposition}[Divisorial ramification is hidden from the \'etale sheet]
\label{prop:etale-sheet-hidden-inertia}
Suppose that \(N/K\) is the normal closure of \(L/K\), the residue
characteristic is zero, and \(X\to Y\) is \'etale.  Then
\[
  \mathcal H(\mathcal D)=\Ram^{(1)}(Z/Y).
\]
Equivalently, every divisorial ramification orbit has a conjugate sheet on
which its inertia acts nontrivially, but every such center lies in the
deleted boundary.
\end{proposition}

\begin{proof}
The inclusion
\(\mathcal H(\mathcal D)\subseteq\Ram^{(1)}(Z/Y)\) is part of the
definition.  Conversely, let \(E\in\Ram^{(1)}(Z/Y)\).  In residue
characteristic zero, \(I_E\neq1\).  Since \(N\) is the normal closure of
\(L/K\), Proposition~\ref{prop:conjugate-embedding-action} gives
\(\core_G(H)=1\), and Proposition~\ref{prop:core-free-sheet-action}
supplies an \(\omega\in\mathcal Q_E\) with \(i_E(\omega)>1\).
For every \(\omega'\in\mathcal Q_E\), Proposition
\ref{prop:visibility-to-boundary} gives
\[
  i_E(\omega')>1
  \quad\Longrightarrow\quad
  \operatorname{cent}(E,\omega')\in\partial X.
\]
These are exactly the two conditions defining
\(E\in\mathcal H(\mathcal D)\).
\end{proof}

\subsection{Automorphism detection}
\label{sec:automorphism-detection}

\begin{theorem}[Automorphism criterion]
\label{thm:automorphism-criterion}
For an arbitrary polynomial self-map
\(F\colon\A^n_{\C}\to\A^n_{\C}\), the following are equivalent:
\begin{enumerate}[label=\textup{(\arabic*)}]
  \item \(F\) is a polynomial automorphism;
  \item \(F\) is injective on \(\A^n(\C)\);
  \item \(I_R(F)=I_\Delta\);
  \item \(\Obs(F)=0\).
\end{enumerate}
\end{theorem}

\begin{proof}
Suppose first that \(F\) has polynomial inverse
\(G=(G_1,\dots,G_n)\).  Then
\[
  x_i-y_i
  =
  G_i(F(x))-G_i(F(y))
  \in I_R
\]
for every \(i\), so \(I_\Delta\subseteq I_R\).  Proposition
\ref{prop:collision-contained-in-diagonal} gives equality.  Thus
\textup{(1)} implies \textup{(3)}.

Suppose \(I_R=I_\Delta\), and let
\(\operatorname{ev}_{(a,b)}\colon S\to\C\) be evaluation at
\((a,b)\).  If \(F(a)=F(b)\), then
\[
  I_R\subseteq\ker(\operatorname{ev}_{(a,b)}).
\]
Therefore \(x_i-y_i\in I_\Delta=I_R\) gives \(a_i-b_i=0\) for every
\(i\).  Thus \(a=b\), proving
\textup{(3)}\(\Rightarrow\)\textup{(2)}.  The
Ax--Grothendieck theorem \cite{Ax1969} gives
\textup{(2)}\(\Rightarrow\)\textup{(1)}.  Finally,
Proposition~\ref{prop:obstruction-kernel} gives
\textup{(3)}\(\Longleftrightarrow\)\textup{(4)}.
\end{proof}

\begin{corollary}[The Jacobian conjecture as obstruction vanishing]
\label{cor:jacobian-conjecture-as-vanishing}
For every \(n\geq1\), the \(n\)-dimensional Jacobian conjecture is
equivalent to each of the statements
\[
  \forall F\colon\A^n_{\C}\to\A^n_{\C},
  \qquad
  \det JF\in\C^\times
  \quad\Longrightarrow\quad
  \Obs(F)=0,
\]
and
\[
  \forall F\colon\A^n_{\C}\to\A^n_{\C},
  \qquad
  \det JF\in\C^\times
  \quad\Longrightarrow\quad
  I_R(F)=I_\Delta.
\]
\end{corollary}

\begin{proof}
Apply Theorem~\ref{thm:automorphism-criterion} to every Keller map.
Its equivalent conditions \textup{(1)}, \textup{(3)}, and \textup{(4)}
give the two displayed formulations.
\end{proof}

\subsection{Generic degree one}

\begin{corollary}[Generic degree-one descent]
\label{cor:generic-degree-one-descent}
If
\[
  L=K
  \qquad\text{and}\qquad
  A\text{ is flat over }B,
\]
then
\[
  \Obs(F)=0.
\]
\end{corollary}

\begin{proof}
When \(L=K\), Theorem~\ref{thm:generic-collision-base-change}
identifies the base-changed diagonal with tensor multiplication:
\[
  \ker(1\otimes\bar\mu_F)
  \cong
  \ker\!\left(
    K\otimes_KK\xrightarrow{\mu_K}K
  \right)
  =
  0.
\]
The fiber-product identification of
Theorem~\ref{thm:collision-tensor-product} is
\[
  C_F\xrightarrow{\sim}A\otimes_BA.
\]
Since \(A\) is flat over \(B\), the tensor product \(A\otimes_BA\), and
therefore \(C_F\), is flat over \(B\).  The localization map
\[
  \iota_C\colon C_F\hookrightarrow K\otimes_BC_F
\]
is therefore injective, and
\[
  \iota_C\bigl(\ker\bar\mu_F\bigr)
  \subseteq
  \ker(1\otimes\bar\mu_F)
  =
  0.
\]
Thus \(\ker\bar\mu_F=0\).  Proposition
\ref{prop:obstruction-kernel} identifies
\[
  \ker(\bar\mu_F)=\Obs(F)=I_\Delta/I_R,
\]
so \(\Obs(F)=0\).
\end{proof}

\begin{leanfiles}
\leanpath{CollisionIdeals/General/Collision/Diagonal.lean},
\leanpath{CollisionIdeals/General/Collision/Ideal.lean},
\leanpath{CollisionIdeals/General/Collision/CollisionDiagonal.lean},
\leanpath{CollisionIdeals/General/FiberProduct/Polynomial.lean},
\leanpath{CollisionIdeals/General/ResidualCollision/Scheme.lean},
\leanpath{CollisionIdeals/General/GenericFiber/FunctionField.lean},
\leanpath{CollisionIdeals/General/GenericFiber/CollisionDiagonal.lean},
\leanpath{CollisionIdeals/General/GenericFiber/MarkedRootDecomposition.lean},
\leanpath{CollisionIdeals/General/Automorphism/Equivalences.lean},
\leanpath{CollisionIdeals/General/Automorphism/GenericDegreeOne.lean},
\leanpath{CollisionIdeals/General/Galois/NormalClosure.lean},
\leanpath{CollisionIdeals/General/Galois/Sheets.lean},
\leanpath{CollisionIdeals/General/Galois/DecompositionSheets.lean},
\leanpath{CollisionIdeals/General/Normalization/Polynomial.lean},
\leanpath{CollisionIdeals/General/Normalization/Diagram.lean}, and
\leanpath{CollisionIdeals/General/Normalization/VisibleRamification.lean}.
The Lean form of
Theorem~\ref{thm:generic-collision-base-change} takes surjectivity of
\(\theta\) as a parameter; the proof above derives it from generic
finiteness.  The normal-closure and group-theoretic sheet objects are
formalized, while the displayed coset--embedding correspondence is not yet
packaged.  Geometric centers and the equality between group-theoretic and
valuation-theoretic ramification indices remain supplied comparison data
through the normalization-diagram and ramification-realization interfaces.
\end{leanfiles}

\section{Planar vanishing}
\label{sec:planar-vanishing}

We now specialize the dimension-independent collision, normalization,
and inertia constructions of Section~\ref{sec:general-construction} to
dimension two.  The purpose of this section is to identify a concrete
boundary obstruction, not to assume it away implicitly: the secant
calculation separates the diagonal from the off-diagonal collision sheet,
while the normalization calculation shows exactly how nontrivial inertia
would have to escape every relevant conjugate affine sheet.

Let
\[
  F=(P,Q)\colon\A^2_{\C}\longrightarrow\A^2_{\C}
\]
be a polynomial map satisfying
\(\det JF=c\in\C^\times\).

Set, from this point onward,
\[
  A=\C[x_1,x_2],
  \qquad
  B=\C[P,Q]\subseteq A,
  \qquad
  K=\C(P,Q),
  \qquad
  L=\C(x_1,x_2),
\]
and let \(X=\Spec A\), \(Y=\Spec B\).

\begin{proposition}[Planar Keller geometry]
\label{prop:planar-keller-local-geometry}
\[
  \C[u,v]\xrightarrow{\sim}B,
  \qquad
  u\longmapsto P,\quad v\longmapsto Q.
\]
Moreover,
\[
  X\xrightarrow{\ f_F\ }Y\ \text{is \'etale},
  \qquad
  A\ \text{is flat over }B,
  \qquad
  [L:K]<\infty
  \ \text{and }L/K\text{ is separable}.
\]
\end{proposition}

\begin{proof}
The identity
\[
  dP\wedge dQ
  =
  \left(
    \frac{\partial P}{\partial x_1}
    \frac{\partial Q}{\partial x_2}
    -
    \frac{\partial P}{\partial x_2}
    \frac{\partial Q}{\partial x_1}
  \right)dx_1\wedge dx_2
  =
  c\,dx_1\wedge dx_2
  \neq0
\]
shows that \(P,Q\) are algebraically independent.  Thus
\(B\cong\C[u,v]\) and \(Y\cong\A^2_{\C}\).  Moreover,
\[
  A
  \cong
  \C[u,v,U,V]\big/\bigl(P(U,V)-u,\ Q(U,V)-v\bigr)
\]
as a \(\C[u,v]\)-algebra, and its relative Jacobian determinant is
\(c\in A^\times\).  The
Jacobian criterion therefore makes \(B\to A\) \'etale.  In particular
it is flat and quasi-finite.  Since \(f_F\) is dominant and
quasi-finite, its generic fiber is zero-dimensional, so
\([L:K]<\infty\).  The morphism remains \'etale at the generic point;
hence the finite extension \(L/K\) is separable.
\end{proof}

\begin{proposition}[Finite planar Keller rigidity]
\label{prop:finite-planar-keller-rigidity}
If the planar Keller morphism
\(f_F\colon\A^2_{\C}\to\A^2_{\C}\) is finite, then \(F\) is a
polynomial automorphism.
\end{proposition}

\begin{proof}
Proposition~\ref{prop:planar-keller-local-geometry} makes \(f_F\) \'etale.
Thus a finite \(f_F\) is a connected finite \'etale cover of
\(\A^2_{\C}\).  Finite-\'etale rigidity of the affine plane,
Theorem~\ref{thm:a2-etale-rigidity}, makes it an isomorphism.  Its inverse
is a morphism of affine schemes and hence is given by polynomials.
\end{proof}

\begin{definition}[Planar generic degree]
\label{def:planar-generic-degree}
The generic degree of \(F\) is
\[
  d_F:=\deg_{\mathrm{gen}}(F):=[L:K].
\]
Because \(L/K\) is finite and separable, \(d_F\) is also the number of
\(K\)-embeddings \(L\hookrightarrow\overline K\), equivalently the number
of geometric points in the generic fiber of \(f_F\).
\end{definition}

\begin{proposition}[Keller inverse-Jacobian frame]
\label{prop:keller-inverse-jacobian-frame}
Define polynomial derivations of \(A=\C[x_1,x_2]\) by
\[
  \partial_P
  :=c^{-1}\left(
      Q_{x_2}\frac{\partial}{\partial x_1}
      -Q_{x_1}\frac{\partial}{\partial x_2}\right),
  \qquad
  \partial_Q
  :=c^{-1}\left(
      -P_{x_2}\frac{\partial}{\partial x_1}
      +P_{x_1}\frac{\partial}{\partial x_2}\right).
\]
Then
\[
  \partial_P(P)=1,
  \quad \partial_P(Q)=0,
  \quad \partial_Q(P)=0,
  \quad \partial_Q(Q)=1.
\]
They extend uniquely to derivations of \(L=\Frac(A)\), and their
restrictions to \(K=\C(P,Q)\) are the coordinate derivations with respect
to \(P,Q\).
\end{proposition}

\begin{proof}
The first four identities are the inverse-matrix identity for \(JF\), using
\(\det JF=c\).  A derivation extends uniquely to a localization, hence to
\(L=\Frac(A)\).  The four identities then give the asserted restrictions
to \(K\).
\end{proof}

\subsection{The planar secant projector}

Let \(\Delta=V(I_\Delta)\subseteq\Spec S\).  For \(a\in S\), write
\(a|_\Delta\) for its image under \(S\to S/I_\Delta\), and apply this
notation entrywise to matrices.
Let
\[
  \mathbf d=
  \begin{pmatrix}
    x_1-y_1\\
    x_2-y_2
  \end{pmatrix},
  \qquad
  \mathbf r=
  \begin{pmatrix}
    P(x)-P(y)\\
    Q(x)-Q(y)
  \end{pmatrix}.
\]
For \(H\in\C[x_1,x_2]\), let
\[
  D_1H
  =
  \frac{H(x_1,x_2)-H(y_1,x_2)}{x_1-y_1},
  \qquad
  D_2H
  =
  \frac{H(y_1,x_2)-H(y_1,y_2)}{x_2-y_2}.
\]
Both quotients belong to \(S\).  Let
\[
  M_F=
  \begin{pmatrix}
    D_1P&D_2P\\
    D_1Q&D_2Q
  \end{pmatrix}.
\]
Then
\[
  \mathbf r=M_F\mathbf d,
  \qquad
  M_F\big|_\Delta=JF.
\]
Let \(\delta_F=\det(M_F)\).

\begin{proposition}[Planar secant annihilation]
\label{prop:secant-annihilator}
The secant determinant satisfies
\[
  \delta_F I_\Delta\subseteq I_R,
  \qquad
  \delta_F\big|_\Delta=c.
\]
\end{proposition}

\begin{proof}
Multiplying \(\mathbf r=M_F\mathbf d\) by
\(\operatorname{adj}(M_F)\) gives
\[
  \delta_F\mathbf d=\operatorname{adj}(M_F)\mathbf r.
\]
\[
  \mathbf r\in I_R^{\oplus2},
  \qquad
  \operatorname{adj}(M_F)\in M_2(S)
  \quad\Longrightarrow\quad
  \operatorname{adj}(M_F)\mathbf r\in I_R^{\oplus2}.
\]
Therefore
\[
  \delta_F\mathbf d\in I_R^{\oplus2},
\]
or equivalently,
\[
  \forall j\in\{1,2\},
  \qquad
  \delta_F(x_j-y_j)\in I_R.
\]
Hence \(\delta_FI_\Delta\subseteq I_R\).  On the diagonal,
\[
  (D_jF_i)|_\Delta=\frac{\partial F_i}{\partial x_j},
\]
so \(M_F|_\Delta=JF\) and
\(\delta_F|_\Delta=\det JF=c\).
\end{proof}

Let \(\bar\delta_F\) be the class of \(\delta_F\) in \(C_F\), and define
\[
  q_F:=1-c^{-1}\bar\delta_F.
\]
Set also
\[
  e_F:=1-q_F=c^{-1}\bar\delta_F.
\]

\begin{proposition}[Secant separability idempotent]
\label{prop:secant-separability-idempotent}
Under \(C_F\cong A\otimes_BA\), the element \(e_F\) satisfies
\[
  e_F^2=e_F,
  \qquad
  \bar\mu_F(e_F)=1,
  \qquad
  (a\otimes1)e_F=(1\otimes a)e_F
  \quad(a\in A).
\]
Thus \(e_F\) is the diagonal separability idempotent written explicitly by
the planar secant determinant.
\end{proposition}

\begin{proof}
Proposition~\ref{prop:secant-annihilator} gives
\(\bar\delta_FJ=0\), while
\(\bar\delta_F-c\in J\).  Hence
\(\bar\delta_F(\bar\delta_F-c)=0\), so
\(e_F^2=e_F\).
Diagonal evaluation gives
\(\bar\mu_F(\bar\delta_F)=c\), so
\(\bar\mu_F(e_F)=1\).  Finally,
\(a\otimes1-1\otimes a\) belongs to
\(J=\ker(\bar\mu_F)\), while \(e_F\) annihilates \(J\); this is exactly the
last displayed identity.
\end{proof}

\begin{theorem}[Planar collision-sheet decomposition]
\label{thm:collision-sheet-decomposition}
The collision scheme decomposes as
\[
  R_F\cong\Delta_{X/Y}(X)\amalg R_F^\circ,
  \qquad
  C_F\cong A\times C_F^\circ.
\]
The off-diagonal ideal satisfies
\[
  I_{\mathrm{off}}
  =
  I_R+(\delta_F)
  =
  I_R:I_\Delta^\infty.
\]
\end{theorem}

\begin{proof}
By Proposition~\ref{prop:obstruction-kernel},
\[
  J
  :=
  \ker(\bar\mu_F)
  =
  I_\Delta/I_R
  =
  \Obs(F).
\]
Proposition~\ref{prop:secant-annihilator} gives
\[
  \bar\delta_FJ=0,
  \qquad
  \bar\delta_F-c\in J.
\]
Theorem~\ref{thm:abstract-off-diagonal-projector}, applied to these two
relations and \(q_F=1-c^{-1}\bar\delta_F\), yields
\[
  q_F^2=q_F,
  \qquad
  J=C_Fq_F,
  \qquad
  \Ann_{C_F}(J)=C_F(1-q_F),
  \qquad
  J=0\Longleftrightarrow q_F=0.
\]
Because \(1-q_F=c^{-1}\bar\delta_F\),
\[
  \frac{I_R:I_\Delta}{I_R}
  =
  \Ann_{C_F}(J)
  =
  C_F(1-q_F)
  =
  C_F\bar\delta_F
  =
  \frac{I_R+(\delta_F)}{I_R}.
\]
Hence
\[
  I_{\mathrm{off}}
  =
  I_R:I_\Delta
  =
  I_R+(\delta_F).
\]

Idempotence gives \(J^m=C_Fq_F=J\) for every \(m\geq1\).  Therefore
\[
  \frac{I_R:I_\Delta^m}{I_R}
  =
  \Ann_{C_F}(J^m)
  =
  \Ann_{C_F}(J)
  =
  \frac{I_R:I_\Delta}{I_R}.
\]
Thus
\[
  I_R:I_\Delta^\infty
  =
  \bigcup_{m\geq1}(I_R:I_\Delta^m)
  =
  I_R:I_\Delta.
\]

The idempotent \(q_F\) also gives a ring isomorphism
\[
  C_F
  \xrightarrow{\ \sim\ }
  C_F(1-q_F)\times C_Fq_F,
  \qquad
  a\longmapsto\bigl(a(1-q_F),aq_F\bigr),
\]
whose inverse is addition.  The two factors satisfy
\[
  C_F(1-q_F)
  \cong
  C_F/C_Fq_F
  =
  C_F/J
  \xrightarrow[\bar\mu_F]{\ \sim\ }
  A
\]
and
\[
  C_Fq_F
  \cong
  C_F/C_F(1-q_F)
  \cong
  S\big/\bigl(I_R+(\delta_F)\bigr)
  =
  C_F^\circ.
\]
Thus \(C_F\cong A\times C_F^\circ\).  Since the spectrum of a product is
the disjoint union of the two spectra, this isomorphism dualizes to
\[
  R_F\cong\Delta_{X/Y}(X)\amalg R_F^\circ.
\]
\end{proof}

\begin{proposition}[Universal moved-collision factorization]
\label{prop:universal-moved-collision-factorization}
Let \(D\) be a domain and let

\[
  \varphi,\psi:A\longrightarrow D
\]

be \(B\)-algebra maps.  If \(\varphi\ne\psi\), then their collision
evaluation
\[
  \operatorname{ev}_{\varphi,\psi}:C_F=A\otimes_BA\longrightarrow D,
  \qquad a\otimes b\longmapsto\varphi(a)\psi(b),
\]
factors uniquely through the off-diagonal algebra:

\[
  C_F\twoheadrightarrow C_F^\circ
  \xrightarrow{\ \exists!\,\overline{\operatorname{ev}}_{\varphi,\psi}\ }D,
  \qquad
  \operatorname{ev}_{\varphi,\psi}
  =\overline{\operatorname{ev}}_{\varphi,\psi}\circ
    \bigl(C_F\twoheadrightarrow C_F^\circ\bigr).
\]

Equivalently, every collision map pair into a domain which moves at least one
of the two planar coordinates annihilates \(I_{\mathrm{off}}\).
\end{proposition}

\begin{proof}
The two maps agree on \(B\), so the induced pair evaluation kills \(I_R\).
Since \(A=\C[x_1,x_2]\) and \(\varphi\ne\psi\), there is an \(i\) with
\(\varphi(x_i)-\psi(x_i)\ne0\).  Evaluating
Proposition~\ref{prop:secant-annihilator} gives

\[
  \varphi\!\times\!\psi(\delta_F)
  \bigl(\varphi(x_i)-\psi(x_i)\bigr)=0.
\]

The second factor is nonzero and \(D\) is a domain, hence
\(\varphi\!\times\!\psi(\delta_F)=0\).  Therefore the evaluation kills

\[
  I_R+(\delta_F)=I_{\mathrm{off}},
\]

and the quotient universal property gives the unique factorization.
\end{proof}

Corollary~\ref{cor:off-diagonal-empty-criterion} states that
\[
  R_F^\circ=\varnothing
  \quad\Longleftrightarrow\quad
  C_F^\circ=0
  \quad\Longleftrightarrow\quad
  \Obs(F)=0
  \quad\Longleftrightarrow\quad
  I_R=I_\Delta.
\]
The preceding theorem identifies
\[
  C_F^\circ
  =
  S\big/\bigl(I_R+(\delta_F)\bigr),
\]
and Theorem~\ref{thm:abstract-off-diagonal-projector} gives
\(\Obs(F)=0\Longleftrightarrow q_F=0\).  Hence
\[
  q_F=0
  \quad\Longleftrightarrow\quad
  I_R+(\delta_F)=S
  \quad\Longleftrightarrow\quad
  R_F^\circ=\varnothing
  \quad\Longleftrightarrow\quad
  I_R=I_\Delta.
\]
In particular, the secant determinant cuts out the off-diagonal factor:
\[
  R_F^\circ
  =
  \Spec\!\bigl(S/(I_R+(\delta_F))\bigr).
\]

The entries of the full secant matrix give more than its determinant.  Write
\[
  M_F=
  \begin{pmatrix}
    m_{11}&m_{12}\\
    m_{21}&m_{22}
  \end{pmatrix},
  \qquad
  d_i=x_i-y_i,
  \qquad
  r_1=P(x)-P(y),\quad r_2=Q(x)-Q(y),
\]
and define
\[
  \Phi_F=
  \begin{pmatrix}
    m_{22}&-m_{21}\\
    -m_{12}&m_{11}\\
    -d_1&-d_2
  \end{pmatrix}.
\]

\begin{proposition}[Secant Hilbert--Burch resolution]
\label{prop:secant-hilbert-burch}
Assume that \(C_F^\circ\ne0\).  The signed maximal minors of \(\Phi_F\)
are \(r_1,r_2,\delta_F\), and there is an exact sequence
\[
  0\longrightarrow S^2
  \xrightarrow{\ \Phi_F\ }S^3
  \xrightarrow{\ (r_1,r_2,\delta_F)\ }S
  \longrightarrow C_F^\circ\longrightarrow0.
\]
Consequently
\[
  \omega_{C_F^\circ}
  \cong\operatorname{coker}(\Phi_F^{\mathsf t})
  \cong C_F^\circ.
\]
The last isomorphism is induced by
\[
  (u,v)\longmapsto-d_2u+d_1v.
\]
\end{proposition}

\begin{proof}
Direct calculation gives
\[
  \det\!\begin{pmatrix}-m_{12}&m_{11}\\-d_1&-d_2\end{pmatrix}=r_1,
  \quad
  -\det\!\begin{pmatrix}m_{22}&-m_{21}\\-d_1&-d_2\end{pmatrix}=r_2,
  \quad
  \det\!\begin{pmatrix}m_{22}&-m_{21}\\-m_{12}&m_{11}\end{pmatrix}
  =\delta_F.
\]
Thus the maximal-minor ideal is
\(I_R+(\delta_F)=I_{\mathrm{off}}\).  The first projection makes
\(C_F=A\otimes_BA\) an \'etale \(A\)-algebra, so its nonzero direct factor
\(C_F^\circ\) is a regular Cohen--Macaulay surface.  Hence
\(I_{\mathrm{off}}\subset S\) has grade two, and Hilbert--Burch gives the
displayed resolution \cite{StacksComplexExactness}.

Dualizing the resolution over the regular four-dimensional ring \(S\) gives
\(\omega_{C_F^\circ}\cong\operatorname{coker}(\Phi_F^{\mathsf t})\).
The row \((-d_2,d_1)\) kills every row of \(\Phi_F\) modulo
\((r_1,r_2,\delta_F)\), so it induces the stated map to \(C_F^\circ\).
Moreover \((d_1,d_2)C_F^\circ=C_F^\circ\): otherwise a prime containing
both would also contain \(\delta_F\), although
\(\delta_F\equiv c\pmod{(d_1,d_2)}\).  The induced map is therefore
surjective.  Both source and target are rank-one locally free modules on the
regular scheme \(\Spec C_F^\circ\), so it is an isomorphism.
\end{proof}

The final trivialization is intrinsic up to a unit.  Identifying this
particular Hilbert--Burch generator with a chosen volume form requires a
separate comparison; no such identification is used below.

\subsection{Galois specialization and conjugate coordinates}

Fix an algebraic closure \(\overline K\supseteq L\), and let
\[
  \begin{aligned}
  N
    &:=
      K\bigl(\sigma(L):
        \sigma\in\Hom_K(L,\overline K)\bigr),
    &K&\subseteq L\subseteq N\subseteq\overline K,\\
  \overline A
    &:=
      \{a\in L:a\text{ is integral over }B\},
  &
  A_N
    &:=
      \{a\in N:a\text{ is integral over }B\},\\
  \overline X&:=\Spec\overline A,
  &
  Z&:=\Spec A_N.
  \end{aligned}
\]
Thus \(N/K\) is the normal closure of \(L/K\).

\begin{proposition}[Planar finite-cover diagram]
\label{prop:planar-normalization-diagram}
\[
  Z\xrightarrow[\mathrm{finite}]{\pi}
  \overline X\xrightarrow[\mathrm{finite}]{\bar\nu}Y,
  \qquad
  \nu=\bar\nu\circ\pi,
\]
and
\[
  X\xhookrightarrow[\mathrm{open}]{j}\overline X,
  \qquad
  \bar\nu\circ j=f_F.
\]
\end{proposition}

\begin{proof}
Since \(B\cong\C[u,v]\) is a finitely generated normal
\(\C\)-algebra, its integral closures in the finite extensions \(L/K\) and
\(N/K\) are finite over \(B\).  The inclusion \(L\subseteq N\) gives
\(\overline A\subseteq A_N\), so
\[
  B\subseteq\overline A\subseteq A_N.
\]
The morphism \(\bar\nu\) is finite because \(\overline A\) is a finite
\(B\)-module.  If \(a_1,\dots,a_r\) generate \(A_N\) over \(B\), then
they also generate \(A_N\) over \(\overline A\), since
\(B\subseteq\overline A\).  Hence \(\pi\) is finite.

For \(a\in\overline A\), integrality over \(B\subseteq A\) and the
normality of \(A=\C[x_1,x_2]\) give \(a\in A\).  Thus
  \[
  B\subseteq\overline A\subseteq A,
  \]
and \(\Frac(\overline A)=L\): clearing denominators in an algebraic
equation over \(K=\Frac(B)\) shows that for every \(\ell\in L\), some
nonzero \(b\in B\) makes \(b\ell\) integral over \(B\).  Thus
\(\ell=(b\ell)/b\in\Frac(\overline A)\).
The inclusion \(\overline A\subseteq A\) therefore induces a dominant
birational morphism \(j\colon X\to\overline X\).

The morphism \(j\) is quasi-finite: because
\(\bar\nu\circ j=f_F\), each fiber of \(j\) is contained in a fiber of
the quasi-finite morphism \(f_F\).  It is separated because it is a
morphism between affine schemes.  Since
\(\overline X\) is normal, Zariski's Main Theorem makes \(j\) an open
immersion.
\end{proof}

One consequence of this finiteness is useful bookkeeping, but should not be
confused with a vanishing theorem.  If \(E=V(\mathfrak p_E)\subset Z\) is a
prime divisor and \(D=V(\mathfrak p_E\cap B)\subset Y\) is its image, then

\[
  A_N/\mathfrak p_E
  \quad\text{is finite over}\quad
  B/(\mathfrak p_E\cap B).
\]

Thus \(E\to D\) is a finite map of affine curves.  Neither coordinate ring
is thereby finite-dimensional over \(\C\), and this curvewise finiteness does
not bound poles in the direction transverse to \(E\).  More precisely, if
\(e(E/D)\) is the DVR ramification index and
\(f(E/D)=[\kappa(E):\kappa(D)]\) is the degree of the induced curve map,
then
\[
  \sum_{E\mid D}e(E/D)f(E/D)=[N:K].
\]
Riemann--Hurwitz or equal-genus arguments for compactifications of
\(E\to D\) constrain \(f(E/D)\); the inertia calculation in this paper
requires control of the distinct transverse index \(e(E/D)\).  This is why
curvewise finiteness alone does not eliminate a ramified divisor.

\subsection{Global and divisorial formulations}
\label{sec:two-planar-finiteness-targets}

Let \(\mathfrak b\subseteq\overline A\) be the radical ideal defining
\(\partial X=\overline X\setminus j(X)\).  Because \(\overline X\) is
affine, local cohomology with support in the boundary is
\[
  H^1_{\partial X}(\overline X,\mathcal O_{\overline X})
  =H^1_{\mathfrak b}(\overline A).
\]

\begin{definition}[Global and divisorial finiteness formulations]
\label{def:two-planar-finiteness-targets}
For a planar Keller map \(F\), define:
\begin{enumerate}
  \item The \emph{global boundary formulation}, PlanarBoundaryCoherence,
  denoted
  \(\operatorname{PBC}(F)\), is the assertion that
  \[
    H^1_{\mathfrak b}(\overline A)
    \quad\text{is a finite (equivalently, finitely generated)
    \(\overline A\)-module}.
  \]
  \item The \emph{divisorial differential formulation},
  PlanarRamificationRigidity, denoted
  \(\operatorname{PRR}(F)\), is the assertion that
  \[
    \Omega_{A_N/B}
    \quad\text{has finite length as an \(A_N\)-module}.
  \]
\end{enumerate}
\end{definition}

Here finiteness does not repeat the finiteness already proved in
Proposition~\ref{prop:planar-normalization-diagram}.  The rings
\(\overline A\) and \(A_N\) are finite \(B\)-modules because they are
integral closures in finite field extensions.  Consequently
\(\overline X\) and \(Z\) are reduced, irreducible affine schemes of finite
type over the algebraically closed field \(\C\), hence affine varieties in
the usual sense.  Ordinary higher cohomology of quasi-coherent sheaves on
these affine schemes vanishes (the affine-cohomology theorem, or the
vanishing half of Serre's affineness criterion)
\cite{StacksAffineCohomology}.  Neither fact implies that
the boundary local cohomology \(H^1_{\mathfrak b}(\overline A)\) is finite:
local cohomology measures sections with poles along the omitted closed set.
Similarly, \(\Omega_{A_N/B}\) is automatically a finite \(A_N\)-module,
because \(A_N\) is a finite-type \(B\)-algebra, but requiring it to have
finite \emph{length} is the stronger assertion that its support has
dimension zero.

\begin{proposition}[The boundary-coherence implication]
\label{prop:planar-boundary-coherence-route}
Planar boundary coherence gives the implication chain
\[
  \operatorname{PBC}(F)
  \Longrightarrow
  \partial X=\varnothing
  \Longrightarrow
  f_F\text{ is finite \'{e}tale}
  \Longrightarrow
  F\text{ is a polynomial automorphism}.
\]
\end{proposition}

\begin{proof}
Because \(\overline A\) is a domain and \(j(X)\) contains its generic point,
\(H^0_{\mathfrak b}(\overline A)=0\).  The local-cohomology sequence for
\(j(X)\subseteq\overline X\), together with affineness of \(\overline X\),
therefore gives
\[
  0\longrightarrow\overline A
  \longrightarrow A
  \longrightarrow H^1_{\mathfrak b}(\overline A)
  \longrightarrow0.
\]
Here \(\Gamma(X,\mathcal O_X)=A\); this is the standard affine
open-complement local-cohomology sequence
\cite{StacksLocalCohomologySequence}.
If the last term is a finite \(\overline A\)-module, then so is \(A\).
Thus the open immersion \(j\) is finite.  Its image is both open and closed;
since \(\overline X\) is integral and \(X\neq\varnothing\), it is all of
\(\overline X\).  Hence \(\partial X=\varnothing\), so
\(f_F=\bar\nu\) is finite.  It is already \'{e}tale by
Proposition~\ref{prop:planar-keller-local-geometry}, and
Proposition~\ref{prop:finite-planar-keller-rigidity} finishes the proof.
\end{proof}

In particular, \(\operatorname{PBC}(F)\) implies
\(\operatorname{PRR}(F)\): after \(F\) is an automorphism one has
\(A_N=B\), so \(\Omega_{A_N/B}=0\).  The converse does not follow by a
direct module-theoretic comparison.  The ramification-rigidity endgame below
nevertheless shows that, for a planar Keller map, \(\operatorname{PRR}(F)\)
also implies automorphy and hence \(\operatorname{PBC}(F)\).  The two
criteria record different mechanisms: ramification rigidity first eliminates
divisorial ramification, whereas boundary coherence eliminates the entire
deleted boundary at once.

Let \(S_F\subseteq Y\) denote the nonproperness set: the complement of
the largest open subset \(U\subseteq Y\) over which
\(f_F^{-1}(U)\to U\) is proper.

\begin{proposition}[The deleted boundary and properness]
\label{prop:planar-boundary-properness}
Set-theoretically,
\[
  S_F=\bar\nu(\partial X).
\]
Consequently, the following are equivalent:
\[
  f_F\text{ is proper},
  \qquad
  S_F=\varnothing,
  \qquad
  \partial X=\varnothing.
\]
Under any of these conditions, \(F\) is a polynomial automorphism.
\end{proposition}

\begin{proof}
The set \(\bar\nu(\partial X)\) is closed because \(\bar\nu\) is finite.
Over
\[
  U=Y\setminus\bar\nu(\partial X),
\]
the open immersion identifies \(f_F^{-1}(U)\) with
\(\bar\nu^{-1}(U)\); hence the restriction of \(f_F\) is finite and
therefore proper.  Thus \(S_F\subseteq\bar\nu(\partial X)\).

Conversely, suppose that \(y\in\bar\nu(\partial X)\) and that \(f_F\)
were proper over a neighborhood of \(y\).  Since \(f_F\) is quasi-finite,
it would be finite after shrinking to an affine neighborhood \(V\) of
\(y\).  The normal finite \(V\)-scheme \(f_F^{-1}(V)\), with function
field \(L\), would then be the normalization of \(V\) in \(L\).  By
uniqueness of normalization it would equal \(\bar\nu^{-1}(V)\), contrary
to the existence of a boundary point above \(y\).  Therefore
\(\bar\nu(\partial X)\subseteq S_F\), proving the equality and the three
equivalences.

If they hold, \(f_F=\bar\nu\) is finite and \'etale.  Since
\(Y\cong\A^2_{\C}\), Theorem~\ref{thm:a2-etale-rigidity} makes this
connected cover trivial, so \(F\) is a polynomial automorphism.
\end{proof}

Thus the deleted normalization boundary is not auxiliary: it is precisely
the geometric source of nonproperness.  For a nonsingular polynomial map of
the complex plane, a nonempty \(S_F\) is known to be a curve with one point
at infinity \cite{ChauNonproperValueSet}.  The inertia calculations below
identify which divisorial part of this nonproperness can support the normal
closure.

Here dimension two materially sharpens the normalization problem.  Both
\(\overline X\) and \(Z\) are normal Noetherian surfaces finite over the
regular surface \(Y\).  A normal Noetherian scheme of dimension at most two
is Cohen--Macaulay \cite{StacksNormalSurfaceCM}; since the finite maps have
zero-dimensional fibers, miracle flatness makes them finite flat
\cite{StacksMiracleFlatness}.  Their coordinate rings are therefore finite
locally free \(B\)-modules, so the trace, different, and ramification divisor
belong to the finite-flat surface setting.

Moreover, \(j(X)\) is a dense affine open in the separated Noetherian surface
\(\overline X\).  Every irreducible component of
\(\partial X=\overline X\setminus j(X)\) consequently has codimension one
\cite{StacksAffineComplement}: a nonempty deleted boundary is a union of
curves.  Purity of the branch locus rules out ramification supported only in
codimension two.  Thus the remaining planar question is the following
concrete one: can a nonzero ramification divisor on \(Z\) have every
inertia-moving conjugate center on the deleted boundary curves?  The
coordinate criterion below gives an exact valuation formulation of this
question.

Characteristic zero also makes the ramification divisor of the intermediate
normalization explicit.  Write
\[
  R_{\bar\nu}
  =
  \sum_{C\in\overline X^{(1)}}
    \bigl(e(C/\bar\nu(C))-1\bigr)C.
\]
The different exponent is \(e-1\) at each divisorial valuation, and the
\'etaleness of \(f_F\) gives
\[
  \operatorname{Supp}(R_{\bar\nu})\subseteq\partial X.
\]
In canonical-divisor language, \(R_{\bar\nu}\) is the relative canonical
divisor of the finite map.  Its vanishing is therefore the finite-cover
analogue of crepancy; this is a ramification statement, not a consequence
of resolving the isolated singularities of the normal surface
\(\overline X\).

\begin{proposition}[Boundary divisor-class rigidity]
\label{prop:boundary-divisor-class-rigidity}
Let \(D\) be an effective Weil divisor on \(\overline X\) supported on
\(\partial X\).  If \([D]\) is torsion in
\(\operatorname{Cl}(\overline X)\), then \(D=0\).
\end{proposition}

\begin{proof}
Choose \(m>0\) and \(h\in L^\times\) with
\[
  \operatorname{div}_{\overline X}(h)=mD.
\]
The divisor on the right is effective.  Since \(\overline X\) is normal,
the height-one valuation criterion for regularity gives
\(h\in\Gamma(\overline X,\mathcal O_{\overline X})=\overline A\).
Because \(D\) is supported on the deleted boundary,
\(\operatorname{div}_{X}(h|_X)=0\).  Since
\(A=\C[x_1,x_2]\) is a unique factorization domain, this gives
\[
  h|_X\in A^\times
  =\C[x_1,x_2]^\times
  =\C^\times.
\]
The inclusion \(\overline A\subseteq A\) is injective, so \(h\) itself is
constant.  Therefore \(mD=0\), and hence \(D=0\).
\end{proof}

\begin{corollary}[Different-class criterion for planar rigidity]
\label{cor:different-class-planar-rigidity}
If \([R_{\bar\nu}]\) is torsion in
\(\operatorname{Cl}(\overline X)\), then \(F\) is a polynomial
automorphism.  In particular, this holds if the different is principal or
if \(\operatorname{Cl}(\overline X)\) is torsion.
\end{corollary}

\begin{proof}
Proposition~\ref{prop:boundary-divisor-class-rigidity} gives
\(R_{\bar\nu}=0\).  Purity of the branch locus then makes the finite map
\(\bar\nu\) \'etale, and finite-\'etale rigidity of \(\A^2_{\C}\) makes it
an isomorphism.  Thus \(L=K\).  Proposition
\ref{prop:planar-keller-local-geometry} gives that \(A\) is flat over
\(B\), so Corollary~\ref{cor:generic-degree-one-descent} gives
\(\Obs(F)=0\).  Theorem~\ref{thm:automorphism-criterion} finishes the
proof.
\end{proof}

The following specialization recovers the classical Galois case of the
Jacobian conjecture \cite{Campbell1973,Razar1979,Wright1981} in the boundary
language of this section.

\begin{theorem}[Boundary-theoretic Galois rigidity]
\label{thm:planar-galois-rigidity}
If the marked function-field extension \(L/K\) is normal, equivalently
Galois, then
\[
  L=K
  \qquad\text{and}\qquad
  F\in\Aut_{\mathrm{poly}}(\A^2_{\C}).
\]
\end{theorem}

\begin{proof}
Proposition~\ref{prop:planar-keller-local-geometry} shows that \(L/K\) is
finite and separable, so normality is equivalent to the Galois condition.
Let \(B_1,\dots,B_r\) be the prime divisors of \(Y\) that are images of
components of \(R_{\bar\nu}\).  Since \(B=\C[P,Q]\) is a unique
factorization domain, choose an irreducible \(b_i\in B\) such that
\[
  B_i=\operatorname{div}_Y(b_i).
\]

The Galois group preserves the integral closure \(\overline X\) and acts
transitively on the prime divisors above a fixed \(B_i\), so their
ramification indices have a common value \(e_i\)
\cite{StacksGaloisIntegralClosure}.  Put
\[
  D_i=
  \sum_{\substack{C\in\overline X^{(1)}\\ \bar\nu(C)=B_i}} C.
\]
For every \(C\mid B_i\), the valuation formula for pullback gives
\[
  \operatorname{ord}_C(b_i)
  =e(C/B_i)\operatorname{ord}_{B_i}(b_i)=e_i.
\]
There are no other height-one zeros of \(b_i\) on \(\overline X\), and
hence
\[
  \operatorname{div}_{\overline X}(b_i)
  =\bar\nu^*B_i=e_iD_i.
\]
These are the standard Galois valuation and divisor-pullback formulas
\cite{StacksDVRExtensions,StacksRamificationTheory}.

The residue characteristic is zero, so the divisorial extensions are tame
and the different exponent at every \(C\mid B_i\) is \(e_i-1\)
\cite{StacksRiemannHurwitzTame}.  Grouping
the components of the ramification divisor by their images therefore gives
\[
  R_{\bar\nu}=\sum_{i=1}^r(e_i-1)D_i.
\]
Set \(m=\prod_{i=1}^r e_i\), with the empty product understood as \(1\),
and define
\[
  h=\prod_{i=1}^r b_i^{\,m(e_i-1)/e_i}\in K^\times\subseteq L^\times.
\]
Every exponent is a nonnegative integer, and the preceding pullback formula
yields
\[
  \operatorname{div}_{\overline X}(h)
  =\sum_{i=1}^r\frac{m(e_i-1)}{e_i}\,e_iD_i
  =mR_{\bar\nu}.
\]
Thus \([R_{\bar\nu}]\) is torsion in
\(\operatorname{Cl}(\overline X)\), and Corollary
\ref{cor:different-class-planar-rigidity} gives the result.
\end{proof}

\begin{corollary}[Quadratic planar rigidity]
\label{cor:quadratic-planar-rigidity}
No planar Keller map has generic degree \(d_F=2\).  Consequently,
\[
  d_F\leq2
  \quad\Longrightarrow\quad
  d_F=1
  \quad\Longrightarrow\quad
  F\in\Aut_{\mathrm{poly}}(\A^2_{\C}).
\]
In particular, every nonautomorphic planar Keller map would have
\[
  d_F\geq3
  \qquad\text{and}\qquad
  L/K\text{ nonnormal}.
\]
\end{corollary}

\begin{proof}
If \(d_F=2\), then the separable quadratic extension \(L/K\) is Galois.
Theorem~\ref{thm:planar-galois-rigidity} would give \(L=K\), contradicting
\([L:K]=2\).  Hence \(d_F\leq2\) forces \(d_F=1\); Corollary
\ref{cor:generic-degree-one-descent} and Theorem
\ref{thm:automorphism-criterion} give automorphy.  Finally, Theorem
\ref{thm:planar-galois-rigidity} rules out normality in every
nonautomorphic case.
\end{proof}

Outside this normal/Galois specialization, the divisor-theoretic problem is
to determine whether the finite-flat trace and different, together with
\(\overline A\subseteq\C[x_1,x_2]\), force
\([R_{\bar\nu}]\) to be torsion.  The different-class criterion proves the
consequence of such torsion; it does not assert it for a general nonnormal
extension.  Corollary~\ref{cor:quadratic-planar-rigidity} shows that only
nonnormal extensions of generic degree at least three remain.

It is important that this argument cannot simply be rerun for the always
Galois extension \(N/K\).  Boundary divisor-class rigidity applies on
\(\overline X\) because it contains the marked affine plane
\(X\simeq\A^2_{\C}\) as a dense open.  If \(L/K\) is nonnormal, the Galois
normalization \(Z\) instead maps finitely to every conjugate intermediate
model and need not contain a marked affine-plane open whose deleted boundary
supports all of \(R_{Z/Y}\).  This is precisely why the pulled-back
conjugate boundaries and moving-sheet condition are needed.

Let
\[
  \begin{gathered}
    \mathcal D_F=
    \left(
      Z\xrightarrow{\pi}\overline X\xrightarrow{\bar\nu}Y,\,
      j\colon X\hookrightarrow\overline X
    \right),
    \qquad
    \partial X=\overline X\setminus j(X),\\
    G=\Gal(N/K),
    \qquad
    H=\Gal(N/L).
  \end{gathered}
\]

The abstract conjugate centers now admit a direct two-coordinate test.
For a representative \(g\in G\) of
\(\omega=D_EgH\), the conjugate affine sheet is generated inside \(N\) by
\[
  A_g=\C[g(x_1),g(x_2)].
\]

The two derivations of Proposition
\ref{prop:keller-inverse-jacobian-frame}, restricted to \(K\), extend
uniquely across the finite separable extension \(N/K\).  Indeed, if
\(z\in N\) has separable minimal polynomial \(f\in K[X]\), any extension
is forced by
\[
  \widetilde\partial(z)
  =-\frac{(\partial f)(z)}{f'(z)},
  \qquad f'(z)\ne0.
\]
For \(g\in G\), both
\(g\widetilde\partial g^{-1}\) and \(\widetilde\partial\) extend the same
derivation of \(K\), so uniqueness makes them equal.  Consequently
\[
  \widetilde\partial(g(a))=g(\partial(a))\in g(A)
  \qquad(a\in A),
\]
and the extended inverse-Jacobian frame preserves every \(A_g\).  No
preservation of the integral order \(T=A_N\) is asserted.

\begin{proposition}[Planar coordinate test for a conjugate center]
\label{prop:planar-conjugate-coordinate-test}
For \(E\in Z^{(1)}\) and \(\omega=D_EgH\in\mathcal Q_E\),
\[
  \operatorname{cent}(E,\omega)\in j(X)
  \quad\Longleftrightarrow\quad
  w_E(g(x_1))\geq0
  \ \text{and}\
  w_E(g(x_2))\geq0.
\]
Consequently,
\[
  \operatorname{cent}(E,\omega)\in\partial X
  \quad\Longleftrightarrow\quad
  \min\{w_E(g(x_1)),w_E(g(x_2))\}<0.
\]
\end{proposition}

\begin{proof}
By definition,
\(u_\omega=(g^{-1}w_E)|_L\), so
\[
  u_\omega(x_i)=w_E(g(x_i))
  \qquad(i=1,2).
\]
The center of \(u_\omega\) lies on the affine open
\(X=\Spec\C[x_1,x_2]\) exactly when its valuation ring contains
\(A\).  Since \(A\) is generated by \(x_1,x_2\), this is equivalent to
the two displayed nonnegativity conditions.  Negating them gives the
boundary criterion.
\end{proof}

The early secant estimate has a precise interpretation on these conjugate
generic sheets.  For \(g\in G\), let
\[
  \delta_{F,g}
  :=
  \delta_F\bigl(x_1,x_2,g(x_1),g(x_2)\bigr)\in N,
  \qquad
  q_{F,g}:=1-c^{-1}\delta_{F,g}.
\]

\begin{proposition}[Secant evaluation on conjugate sheets]
\label{prop:secant-conjugate-evaluation}
For every \(g\in G\),
\[
  \delta_{F,g}\bigl(x_i-g(x_i)\bigr)=0
  \qquad(i=1,2).
\]
Consequently,
\[
  \delta_{F,g}
  =
  \begin{cases}
    c,&g\in H,\\
    0,&g\notin H,
  \end{cases}
  \qquad
  q_{F,g}
  =
  \begin{cases}
    0,&g\in H,\\
    1,&g\notin H.
  \end{cases}
\]
Thus the planar secant projector is exactly the characteristic function of
the nonmarked generic conjugate sheets.
\end{proposition}

\begin{proof}
The marked embedding \(L\hookrightarrow N\) and its \(g\)-conjugate agree
on \(K\), so together they define an evaluation map from the generic
collision algebra \(L\otimes_KL\) to \(N\).  Evaluating
Proposition~\ref{prop:secant-annihilator} under this map gives the first
display.  If \(g\in H=\Gal(N/L)\), then \(g(x_i)=x_i\) for both coordinates,
and diagonal evaluation gives \(\delta_{F,g}=c\).  If \(g\notin H\), then
some \(x_i\) is moved, since \(L=\C(x_1,x_2)\).  The corresponding factor
\(x_i-g(x_i)\) is nonzero in the field \(N\), so the first display forces
\(\delta_{F,g}=0\).  The formulas for \(q_{F,g}\) follow.
\end{proof}

The complementary idempotent \(e_F=1-q_F\) gives an explicit
trace-theoretic bridge.  Under the standard finite Galois decomposition
\[
  N\otimes_KN\xrightarrow{\sim}\prod_{\sigma\in G}N,
  \qquad
  a\otimes b\longmapsto\bigl(a\,\sigma(b)\bigr)_{\sigma},
\]
the generic base change of \(e_F\) is the characteristic function of
\(H\).  Let \(e_g=(g\otimes g)e_F\); then \(e_g\) is the characteristic
function of \(gHg^{-1}\).

\begin{proposition}[Relative coevaluation on a moved sheet]
\label{prop:relative-coevaluation-moved-sheet}
Let \(C\leq G\) have prime order and suppose
\(C\nsubseteq gHg^{-1}\).  Under the relative Galois decomposition
\[
  N\otimes_{N^C}N\xrightarrow{\sim}\prod_{\sigma\in C}N,
  \qquad
  a\otimes b\longmapsto\bigl(a\,\sigma(b)\bigr)_{\sigma},
\]
let
\[
  \rho_C:N\otimes_KN\longrightarrow N\otimes_{N^C}N
\]
be the canonical quotient.  Then
\(\rho_C(e_g)=\mathbf1_{\{1\}}\).  Consequently \(\rho_C(e_g)\) is the
trace coevaluation tensor for \(N/N^C\).  In particular, if its image is
written as a finite sum
\[
  \rho_C(e_g)=\sum_\nu a_{\nu,g}\otimes b_{\nu,g}
  \quad\text{in }N\otimes_{N^C}N,
\]
then every \(z\in N\) satisfies the explicit reconstruction formula
\[
  z=\sum_\nu a_{\nu,g}\,
       \operatorname{Tr}_{N/N^C}(b_{\nu,g}z).
\]
\end{proposition}

\begin{proof}
Restricting the characteristic function of \(gHg^{-1}\) from \(G\) to
\(C\) gives the characteristic function of
\(C\cap gHg^{-1}\).  Since \(C\) has prime order and is not contained in
\(gHg^{-1}\), this intersection is trivial.  The relative trace-pairing
identification sends \(\mathbf1_{\{1\}}\) to \(\operatorname{id}_N\),
which gives the reconstruction formula.
\end{proof}

Proposition
\ref{prop:conjugate-embedding-action} gives
\[
  e_N:=\prod_{g\in G}e_g
  =\mathbf1_{\core_G(H)}
  =\mathbf1_{\{1\}}.
\]
Finally, the trace pairing gives a canonical \(K\)-linear isomorphism
\[
  \Theta_N:N\otimes_KN\xrightarrow{\sim}\operatorname{End}_K(N),
  \qquad
  \Theta_N(a\otimes b)(z)=a\operatorname{Tr}_{N/K}(bz),
\]
and \(\Theta_N(e_N)=\operatorname{id}_N\).  Thus the product of the
conjugate planar secant idempotents is the trace coevaluation tensor for
\(N/K\).  Notice that \(\Theta_N\) is a linear, not a ring, isomorphism.
For the finite locally free order \(T=A_N\), the integral coevaluation tensor
lies in \(T\otimes_B\operatorname{Hom}_B(T,B)\); its generic fiber is
\(e_N\) after the trace identification.  This identifies the finite
trace-dual lattice canonically, but does not yet bound poles of moving-sheet
principal parts.

This follows the exponent-one secant bound completely into the normal
closure, but also locates its limit.  On a nonmarked sheet
\(\delta_{F,g}\) is the zero function; it is therefore a generic-sheet
projector, not a nonzero boundary parameter whose fixed power bounds pole
order.  Any use of the secant matrix for uniform boundary control therefore
needs a further transport statement, using more than the determinant alone.
That possible mechanism is isolated after the completed endgame.

The generic idempotents and the conjugate affine opens play different
roles.  The former separate components of the generic collision algebra,
whereas \(U_g\subset Z\) tests whether the corresponding conjugate center
remains integral.  At a ramified divisor those generic components can meet
in an integral model, so no extension of \(e_g\) as an integral idempotent is
asserted.  The research mechanism below instead attaches to this generic
sheet isolation a fixed--moving boundary contribution supported on

\[
  V\!\left(\mathfrak a_C+\mathfrak j_C^{\mathrm{mov}}\right).
\]

For \(E\in\Ram^{(1)}(Z/Y)\) and
\(\omega\in\mathcal Q_E\), Proposition
\ref{prop:intermediate-ramification-index} gives
\[
  e\bigl(\operatorname{cent}(E,\omega)/\nu(E)\bigr)=i_E(\omega).
\]
Since \(X\to Y\) is \'etale by
Proposition~\ref{prop:planar-keller-local-geometry},
Proposition~\ref{prop:visibility-to-boundary} gives
\[
  i_E(\omega)>1
  \quad\Longrightarrow\quad
  \operatorname{cent}(E,\omega)\in\partial X.
\]
Thus Definition~\ref{def:no-hidden-inertia} specializes to
\[
  \mathcal H(\mathcal D_F)
  =
  \left\{
  E\in\Ram^{(1)}(Z/Y):
  \begin{array}{l}
    \exists\omega\in\mathcal Q_E,\ i_E(\omega)>1,\\
    \forall\omega\in\mathcal Q_E,\quad
      i_E(\omega)>1
      \Longrightarrow
      \operatorname{cent}(E,\omega)\in\partial X
  \end{array}
  \right\}.
\]

Proposition~\ref{prop:planar-conjugate-coordinate-test} makes this locus
fully explicit:
\[
  \mathcal H(\mathcal D_F)
  =
  \left\{
  E\in\Ram^{(1)}(Z/Y):
  \begin{array}{l}
    \forall g\in G,\\[-2pt]
    [I_E:I_E\cap gHg^{-1}]>1
    \Longrightarrow
    \min\{w_E(g(x_1)),w_E(g(x_2))\}<0
  \end{array}
  \right\}.
\]
The existence of at least one \(g\) with nontrivial relative index is
automatic from \(I_E\neq1\) and \(\core_G(H)=1\).  Thus hidden inertia is
not an additional abstract object in the planar proof: it is precisely the
possibility that every inertia-moving conjugate coordinate pair acquires a
pole at the deleted boundary.

\begin{definition}[Planar conjugate coverage]
\label{def:planar-conjugate-coverage}
Write \(\operatorname{Cov}_2(F)\) for the following condition: every
\(E\in\Ram^{(1)}(Z/Y)\) admits a \(g\in G\) such that
\[
  [I_E:I_E\cap gHg^{-1}]>1,
  \qquad
  w_E(g(x_1))\geq0,
  \qquad
  w_E(g(x_2))\geq0.
\]
Thus some conjugate sheet moved by inertia still has an affine center.
\end{definition}

We now construct the canonical boundaries directly on the common Galois
normalization.  For \(g\in G\), set
\[
  \overline A_g=g(\overline A),
  \qquad
  \overline X_g=\Spec\overline A_g,
  \qquad
  X_g=\Spec A_g.
\]
The conjugate of \(j\) is an open immersion
\(X_g\hookrightarrow\overline X_g\), and the inclusion
\(\overline A_g\subseteq A_N\) gives a finite map
\(\pi_g\colon Z\to\overline X_g\).  Let
\[
  U_g:=\pi_g^{-1}(X_g)\subseteq Z,
  \qquad
  D_g:=Z\setminus U_g,
  \qquad
  \mathfrak j_g:=
    I_{A_N}(D_g)
    =
    \bigcap_{\mathfrak p\in D_g}\mathfrak p.
\]
Thus \(\mathfrak j_g\) is canonically the radical vanishing ideal of the
pulled-back deleted boundary and \(D_g=V(\mathfrak j_g)\).  The
construction depends only on \(gH\), since right multiplication by \(H\)
does not change the conjugate embedding of \(L\).

\begin{proposition}[Pulled-back boundary dictionary]
\label{prop:pulled-back-boundary-dictionary}
For every height-one point \(E\) of \(Z\) and every \(g\in G\),
\[
  \eta_E\in U_g
  \quad\Longleftrightarrow\quad
  \mathfrak j_g\nsubseteq\mathfrak p_E
  \quad\Longleftrightarrow\quad
  w_E(g(x_1))\geq0
  \quad\text{and}\quad
  w_E(g(x_2))\geq0.
\]
Equivalently,
\[
  \mathfrak j_g\subseteq\mathfrak p_E
  \quad\Longleftrightarrow\quad
  \eta_E\in D_g
  \quad\Longleftrightarrow\quad
  \min\{w_E(g(x_1)),w_E(g(x_2))\}<0.
\]
\end{proposition}

\begin{proof}
Because \(D_g=V(\mathfrak j_g)\), membership
\(\eta_E\in D_g\) is exactly the containment
\(\mathfrak j_g\subseteq\mathfrak p_E\).  The point
\(\eta_E\) belongs to \(U_g\) precisely when its image under
\(\pi_g\) is centered on \(X_g=\Spec\C[g(x_1),g(x_2)]\).  The valuation
criterion for a center on an affine model says exactly that both conjugate
coordinates have nonnegative valuation.  Taking complements gives the
second display.
\end{proof}

For a nontrivial subgroup \(I\leq G\), define its moving-boundary ideal by
\[
  \mathfrak j_I^{\mathrm{mov}}
  :=
  \sum_{\substack{g\in G\\I\nsubseteq gHg^{-1}}}\mathfrak j_g.
\]
The sum is taken in the common Galois-normalization ring \(A_N\); this is
why the Galois closure enters only at this final specialization.

\begin{proposition}[Moving-sheet coverage dictionary]
\label{prop:moving-sheet-coverage-dictionary}
For \(E\in\Ram^{(1)}(Z/Y)\), one has
\[
  \exists g\in G,\quad
  I_E\nsubseteq gHg^{-1},\quad \eta_E\in U_g
  \quad\Longleftrightarrow\quad
  \mathfrak j_{I_E}^{\mathrm{mov}}\nsubseteq\mathfrak p_E.
\]
Thus eliminating \(E\) is exactly the failure of its moving-boundary
ideal to be contained in its height-one prime.
\end{proposition}

\begin{proof}
A prime contains a sum of ideals exactly when it contains every summand.
Apply Proposition~\ref{prop:pulled-back-boundary-dictionary} to each
inertia-moving \(g\).
\end{proof}

The same common ring gives a finite-group formulation which does not first
choose a ramified divisor.  For a subgroup \(C\leq G\), let
\[
  \mathfrak a_C
  :=
  \bigl(\sigma(t)-t:\sigma\in C,\ t\in A_N\bigr)
  \subseteq A_N.
\]
This is the scheme-theoretic fixed-locus ideal for the action of \(C\) on
\(Z\).  If \(\mathfrak p_E\in\Spec A_N\), then
\[
  \mathfrak a_C\subseteq\mathfrak p_E
  \quad\Longleftrightarrow\quad
  C\leq I_E.
\]
Indeed, the containment says that every element of \(C\) stabilizes
\(\mathfrak p_E\) and acts trivially on its residue field, which is the
definition of inertia at \(E\).

\begin{definition}[Planar boundary separation]
\label{def:planar-boundary-separation}
Write \(\operatorname{Sep}_2(F)\) for the following condition: for every
prime-order subgroup \(C\leq G\),
\[
  \Spec^{(1)}(A_N)
  \cap
  V\!\left(\mathfrak a_C+\mathfrak j_C^{\mathrm{mov}}\right)
  =\varnothing.
\]
Equivalently,
\[
  \forall\mathfrak p\in\Spec^{(1)}(A_N),\qquad
  \mathfrak a_C\subseteq\mathfrak p
  \quad\Longrightarrow\quad
  \mathfrak j_C^{\mathrm{mov}}\nsubseteq\mathfrak p.
\]
In codimension language, every fixed--moving boundary locus
\(V(\mathfrak a_C+\mathfrak j_C^{\mathrm{mov}})\) has codimension at least
two in the normal surface \(Z\).  No assertion is made about isolated
fixed boundary points.
\end{definition}

\begin{proposition}[Boundary separation is the divisorial criterion]
\label{prop:boundary-separation-divisorial-criterion}
For a planar Keller map,
\[
  \operatorname{Sep}_2(F)
  \quad\Longleftrightarrow\quad
  \Ram^{(1)}(Z/Y)=\varnothing.
\]
\end{proposition}

\begin{proof}
Assume boundary separation and suppose that
\(E\in\Ram^{(1)}(Z/Y)\).  Choose a subgroup
\(C\leq I_E\) of prime order; such a subgroup exists by Cauchy's theorem
for the nontrivial finite group \(I_E\).  Then
\(\mathfrak a_C\subseteq\mathfrak p_E\).  If
\(C\nsubseteq gHg^{-1}\), then
\(I_E\nsubseteq gHg^{-1}\), so the relative inertia index on the sheet
represented by \(gH\) is nontrivial.  Proposition
\ref{prop:visibility-to-boundary} places its center in the deleted boundary,
and hence \(\mathfrak j_g\subseteq\mathfrak p_E\).  This holds for every
sheet moved by \(C\), so
\[
  \mathfrak a_C+\mathfrak j_C^{\mathrm{mov}}
  \subseteq\mathfrak p_E,
\]
contrary to boundary separation.

Conversely, if boundary separation fails, there are a prime-order
\(C\leq G\) and a height-one prime \(\mathfrak p_E\) containing
\(\mathfrak a_C+\mathfrak j_C^{\mathrm{mov}}\).  The fixed-locus
containment gives \(C\leq I_E\).  Since \(C\neq1\), the point \(E\) is
ramified, contradicting
\(\Ram^{(1)}(Z/Y)=\varnothing\).
\end{proof}

\begin{proposition}[Coverage eliminates a ramified divisor]
\label{prop:coverage-eliminates-ramified-divisor}
Let \(E\in\Ram^{(1)}(Z/Y)\).  If there is a \(g\in G\) such that
\[
  I_E\nsubseteq gHg^{-1}
  \qquad\text{and}\qquad
  \eta_E\in U_g,
\]
then a contradiction follows.  Consequently,
\[
  \operatorname{Cov}_2(F)
  \quad\Longrightarrow\quad
  \Ram^{(1)}(Z/Y)=\varnothing.
\]
\end{proposition}

\begin{proof}
The subgroup condition gives
\[
  i_E(gH)=[I_E:I_E\cap gHg^{-1}]>1.
\]
On the other hand, \(\eta_E\in U_g\) places the conjugate center on the
affine sheet \(X_g\).  The map \(X_g\to Y\) is the Galois conjugate of
the planar Keller map and is therefore \'etale.  Its local ramification
index at that center is one.  Hence \(1<i_E(gH)=1\), a contradiction.
\end{proof}

This is the point at which a height-one divisor is eliminated.  Purity is
not used in Proposition~\ref{prop:coverage-eliminates-ramified-divisor};
it enters only after the proposition has been applied to every ramified
height-one point.

\subsection{Moving-sheet criteria for planar ramification rigidity}
\label{sec:planar-no-hidden-inertia}

\begin{corollary}[Planar hidden-inertia identification]
\label{cor:planar-codimension-one-unramified}
\[
  \mathcal H(\mathcal D_F)=\Ram^{(1)}(Z/Y).
\]
In particular,
\[
  \mathcal H(\mathcal D_F)=\varnothing
  \quad\Longleftrightarrow\quad
  \Ram^{(1)}(Z/Y)=\varnothing.
\]
These conditions are also equivalent to planar conjugate coverage.
\end{corollary}

\begin{proof}
The equality is Proposition~\ref{prop:etale-sheet-hidden-inertia}, applied
to the planar finite-cover diagram.  Conjugate coverage
implies an empty ramification locus by Proposition
\ref{prop:coverage-eliminates-ramified-divisor}, and hence an empty
hidden-inertia locus by the displayed equality.  Conversely, if the
ramification locus is empty, conjugate coverage holds vacuously.
\end{proof}

\begin{proposition}[The ramification-rigidity dictionary]
\label{prop:planar-ramification-rigidity-dictionary}
For a planar Keller map,
\[
  \operatorname{PRR}(F)
  \quad\Longleftrightarrow\quad
  \Ram^{(1)}(Z/Y)=\varnothing
  \quad\Longleftrightarrow\quad
  \operatorname{Cov}_2(F)
  \quad\Longleftrightarrow\quad
  \operatorname{Sep}_2(F).
\]
Thus moving-sheet coverage and boundary separation are two exact criteria
for PlanarRamificationRigidity, not additional global hypotheses.
\end{proposition}

\begin{proof}
The \(A_N\)-module \(\Omega_{A_N/B}\) is finite, and its generic
localization vanishes because \(N/K\) is separable.  Since \(A_N\) is a
two-dimensional affine \(\C\)-domain, every positive-dimensional closed
subset has a height-one generic point.  Thus this finite module has finite
length exactly when it has no height-one point in its support.  At a height-one
point \(E\), localizing \(B\) first at the image of \(E\) and then at
\(\mathfrak p_E\) gives the relevant finite local extension; the source
local ring is a discrete valuation ring.  In characteristic zero its
residue extension is separable, so
\[
  \Omega_{(A_N)_{\mathfrak p_E}/B_{\mathfrak p_E\cap B}}=0
  \quad\Longleftrightarrow\quad
  I_E=1.
\]
Indeed, vanishing of the localized differential module is the
finite-type unramifiedness criterion, and for a finite extension of
discrete valuation rings in characteristic zero unramifiedness is
equivalent to trivial inertia
\cite{StacksKaehlerDifferent,StacksDVRExtensions,StacksRamificationTheory}.
Consequently finite length is equivalent to the absence of divisorial
ramification.  Proposition~\ref{prop:coverage-eliminates-ramified-divisor}
and vacuity in the unramified case give the equivalence with coverage; the
equivalence with separation is
Proposition~\ref{prop:boundary-separation-divisorial-criterion}.
\end{proof}

Equivalently, the premise
\(\mathcal H(\mathcal D_F)=\varnothing\) says that if a ramified divisor
existed, some inertia-moving conjugate pair \(g(x_1),g(x_2)\) would have
to remain integral and hence centered on an affine sheet.  Planar
\'etaleness says that the same pair must lie at the boundary.  This is the
boundary contradiction used below; the premise supplies its exclusion,
while the preceding results identify exactly what is being excluded.

\subsection{The divisorial endgame}

\begin{theorem}[Purity of the branch locus, specialized]
\label{thm:branch-purity}
Let \(Z\to\A^2_{\C}\) be a finite dominant generically separable
morphism with \(Z\) integral and normal.  If
\(\Ram^{(1)}(Z/\A^2_{\C})=\varnothing\), then the morphism is \'etale.
\end{theorem}

\begin{proof}
At every height-one point of \(Z\), the absence of divisorial
ramification gives ramification index one.  The corresponding residue
field extension is finite and, in characteristic zero, separable.
Hence the morphism is unramified at every height-one point.  Purity of
the branch locus now implies that it is \'etale everywhere
\cite{StacksPurity}.
\end{proof}

\begin{theorem}[Finite-\'etale rigidity of the affine plane]
\label{thm:a2-etale-rigidity}
Every connected finite \'etale cover of \(\A^2_{\C}\) is trivial
\cite[Expos\'e XII, Th\'eor\`eme 5.1]{SGA1}.
\end{theorem}

\begin{corollary}[Triviality of the planar normal closure]
\label{cor:planar-normal-closure-trivial}
If
\[
  \Ram^{(1)}(Z/Y)=\varnothing,
\]
then
\[
  N=K
  \qquad\text{and hence}\qquad
  L=K.
\]
Numerically,
\[
  [L:K]=[N:L]=[N:K]=1.
\]
\end{corollary}

\begin{proof}
By Proposition~\ref{prop:planar-keller-local-geometry},
\(Y\cong\A^2_{\C}\).  Since \(N/K\) is finite Galois,
\(\nu\colon Z\to Y\) is generically separable; hence
Theorem~\ref{thm:branch-purity} makes it finite \'etale.  The inclusion
\[
  A_N\subset N\text{ is a domain},
  \qquad\text{so}\qquad
  Z=\Spec A_N\text{ is connected}.
\]
Theorem~\ref{thm:a2-etale-rigidity} therefore makes
\(\nu\colon Z\to Y\) an isomorphism.  Taking function fields gives
\(N=K\).  Finally,
\(K\subseteq L\subseteq N\) implies \(L=K\).
\end{proof}

\begin{corollary}[The divisorial endgame]
\label{cor:planar-ramification-rigidity-route}
The absence of divisorial ramification gives the implication chain
\[
  \Ram^{(1)}(Z/Y)=\varnothing
  \Longrightarrow
  Z\longrightarrow Y\text{ is finite \'{e}tale}
  \Longrightarrow
  N=K
  \Longrightarrow
  L=K
  \Longrightarrow
  q_F=0
  \Longrightarrow
  I_R=I_\Delta.
\]
\end{corollary}

\begin{proof}
Purity, Theorem~\ref{thm:branch-purity}, makes \(Z\to Y\) finite \'{e}tale,
and
Corollary~\ref{cor:planar-normal-closure-trivial} gives \(N=K\) and
\(L=K\).  Thus \(d_F=1\).  Flatness of \(A\) over \(B\) and
Corollary~\ref{cor:generic-degree-one-descent} give \(\Obs(F)=0\).
The abstract off-diagonal projector and collision criteria then give
\(q_F=0\) and \(I_R=I_\Delta\).
\end{proof}

By Proposition~\ref{prop:planar-ramification-rigidity-dictionary}, any of
\(\operatorname{PRR}(F)\), \(\operatorname{Sep}_2(F)\), or
\(\operatorname{Cov}_2(F)\) supplies the first condition in this single
endgame.  They are respectively differential, ideal-theoretic, and
valuation-theoretic descriptions of its codimension-one premise, rather
than three additional routes.

\subsection{Vanishing of the off-diagonal factor}

The next theorem records all named formulations as exact translations of
the preceding divisorial endgame.

\begin{theorem}[Equivalent boundary criteria for planar vanishing]
\label{thm:planar-vanishing}
For every polynomial map
\(F=(P,Q)\colon\A^2_{\C}\to\A^2_{\C}\) satisfying
\(\det JF=c\in\C^\times\), the following conditions are equivalent:
\[
  \begin{aligned}
    \operatorname{PBC}(F)
    &\quad\Longleftrightarrow\quad \operatorname{PRR}(F)
     \quad\Longleftrightarrow\quad \operatorname{Sep}_2(F)
     \quad\Longleftrightarrow\quad \operatorname{Cov}_2(F)\\
    &\quad\Longleftrightarrow\quad \mathcal H(\mathcal D_F)=\varnothing
     \quad\Longleftrightarrow\quad \Ram^{(1)}(Z/Y)=\varnothing,
  \end{aligned}
\]
and
\[
  \begin{gathered}
    \Ram^{(1)}(Z/Y)=\varnothing
    \quad\Longleftrightarrow\quad
    d_F=1
    \quad\Longleftrightarrow\quad
    \Obs(F)=0,\\
    \Obs(F)=0
    \quad\Longleftrightarrow\quad
    q_F=0
    \quad\Longleftrightarrow\quad
    R_F^\circ=\varnothing
    \quad\Longleftrightarrow\quad
    I_R=I_\Delta.
  \end{gathered}
\]
\end{theorem}

\begin{proof}
Proposition~\ref{prop:planar-boundary-coherence-route} shows that
\(\operatorname{PBC}(F)\) implies that \(F\) is an automorphism.  Then
\(B=A\), \(K=L=N\), \(\partial X=\varnothing\), and
\(\Omega_{A_N/B}=0\), so \(\operatorname{PRR}(F)\) holds.  Conversely,
the argument below shows that \(\operatorname{PRR}(F)\) implies
automorphy; for an automorphism \(\partial X=\varnothing\), and hence
\(H^1_{\partial X}(\overline X,\mathcal O_{\overline X})=0\).  Thus
\(\operatorname{PBC}(F)\) holds as well.

Proposition~\ref{prop:planar-ramification-rigidity-dictionary} identifies
\(\operatorname{PRR}(F)\) and \(\operatorname{Sep}_2(F)\) with the absence
of divisorial ramification.
Proposition~\ref{prop:coverage-eliminates-ramified-divisor} shows directly
that \(\operatorname{Cov}_2(F)\) implies
\(\Ram^{(1)}(Z/Y)=\varnothing\); conversely, if the ramification locus is
empty, coverage holds vacuously.  Corollary
\ref{cor:planar-codimension-one-unramified} identifies these conditions with
\(\mathcal H(\mathcal D_F)=\varnothing\).

Under these conditions, Corollary
\ref{cor:planar-normal-closure-trivial} gives \(N=L=K\), hence \(d_F=1\),
while Proposition
\ref{prop:planar-keller-local-geometry} makes \(A\) flat over \(B\).
Corollary~\ref{cor:generic-degree-one-descent} therefore gives
\(\Obs(F)=0\).  Conversely, if \(d_F=1\), then \(L=K\); flatness and
Corollary~\ref{cor:generic-degree-one-descent} again give \(\Obs(F)=0\).
If \(\Obs(F)=0\), then \(I_R=I_\Delta\), so
Theorem~\ref{thm:automorphism-criterion} makes \(F\) a polynomial
automorphism.  Hence \(B=A\), \(K=L=N\), and
\(\Ram^{(1)}(Z/Y)=\varnothing\).

Corollary~\ref{cor:off-diagonal-empty-criterion} and
Theorem~\ref{thm:abstract-off-diagonal-projector} state
\[
  \begin{gathered}
    R_F^\circ=\varnothing
    \quad\Longleftrightarrow\quad
    C_F^\circ=0
    \quad\Longleftrightarrow\quad
    \Obs(F)=0
    \quad\Longleftrightarrow\quad
    I_R=I_\Delta,\\
    \Obs(F)=0
    \quad\Longleftrightarrow\quad
    q_F=0.
  \end{gathered}
\]
This proves the remaining equivalences.
\end{proof}

\begin{corollary}[Boundary witness for failure of coverage]
\label{cor:planar-boundary-witness}
The equivalent conditions of Theorem
\ref{thm:planar-vanishing} fail if and only if there is an
\(E\in\Ram^{(1)}(Z/Y)\) such that
\[
  \forall g\in G,\qquad
  [I_E:I_E\cap gHg^{-1}]>1
  \quad\Longrightarrow\quad
  \min\{w_E(g(x_1)),w_E(g(x_2))\}<0.
\]
For such an \(E\), at least one inertia-moving \(g\) exists.  Thus failure of
planar collision vanishing is witnessed by a ramified divisor whose
inertia-moving conjugate centers all lie on the deleted boundary.  By
Theorem~\ref{thm:planar-galois-rigidity} and Corollary
\ref{cor:quadratic-planar-rigidity}, its marked extension \(L/K\) is
nonnormal and its generic degree satisfies \(d_F\geq3\).
\end{corollary}

\begin{proof}
If the equivalent conditions fail, then conjugate coverage fails.  Negating
Definition~\ref{def:planar-conjugate-coverage} supplies an
\(E\in\Ram^{(1)}(Z/Y)\) for which every \(g\) of nontrivial relative index
fails at least one valuation inequality, giving the displayed condition.
Conversely, the displayed condition is incompatible with a coverage witness
for \(E\), so conjugate coverage fails.  Core-freeness guarantees the
existence of at least one \(g\) of nontrivial relative index.
\end{proof}

\begin{corollary}[Exact boundary forms of the planar Jacobian conjecture]
\label{cor:planar-jc2-boundary-forms}
The following universal statements are equivalent:
\begin{enumerate}
  \item The classical complex two-dimensional Jacobian conjecture holds:
  every polynomial self-map of \(\A^2_{\C}\) with constant nonzero
  Jacobian determinant is a polynomial automorphism.
  \item Every planar Keller map satisfies \(\operatorname{PBC}(F)\).
  \item Every planar Keller map satisfies \(\operatorname{PRR}(F)\).
  \item Every planar Keller map satisfies \(\operatorname{Sep}_2(F)\).
  \item Every planar Keller map satisfies \(\operatorname{Cov}_2(F)\).
\end{enumerate}
Equivalently, condition \textup{(4)} says that for every planar Keller map
and every prime-order subgroup \(C\leq\operatorname{Gal}(N/K)\),
  \[
    \operatorname{codim}_Z
    V\!\left(\mathfrak a_C+\mathfrak j_C^{\mathrm{mov}}\right)
    \geq2.
  \]
\end{corollary}

\begin{proof}
If \(JC(2)\) holds, every planar Keller map is an automorphism.  Then
\(B=A\), \(K=L=N\), \(\overline X=X\), and \(\partial X=\varnothing\),
so both \(H^1_{\partial X}(\overline X,\mathcal O_{\overline X})\) and
\(\Omega_{A_N/B}\) vanish.  Thus \textup{(1)} implies
\textup{(2)} and \textup{(3)}.  Conversely, Proposition
\ref{prop:planar-boundary-coherence-route} shows that \textup{(2)} implies
\textup{(1)}.  Theorem~\ref{thm:planar-vanishing} identifies, for every
planar Keller map,
\[
  \operatorname{PRR}(F)
  \Longleftrightarrow
  \operatorname{Sep}_2(F)
  \Longleftrightarrow
  \operatorname{Cov}_2(F)
  \Longleftrightarrow
  \Obs(F)=0.
\]
Corollary~\ref{cor:jacobian-conjecture-as-vanishing} identifies universal
obstruction vanishing with classical \(JC(2)\), proving the remaining
equivalences.  The prime-order ideal formulation is
Proposition~\ref{prop:boundary-separation-divisorial-criterion}.
\end{proof}

The planar section therefore gives an exact codimension-one reformulation.
Here ``divisorially coprime'' means that
\(\mathfrak a_C\) and \(\mathfrak j_C^{\mathrm{mov}}\) are contained in no
common height-one prime; it does not mean that their sum is the unit ideal.
Because \(Z\) is an integral surface, the relevant locus has only two
possibilities: it has a codimension-one component, which is precisely the
bad ramified-divisor case, or it has codimension at least two, in which case
only isolated closed points may remain (or the locus is empty).  There is no
nonempty codimension-three case on \(Z\), and purity is why these possible
codimension-two points need not be eliminated separately.

In the good case, conjugate coverage excludes complete boundary-supported
inertia orbits and the collision relation collapses to the diagonal.  In
the bad case, Corollary~\ref{cor:planar-boundary-witness} produces exactly
such a divisorial boundary witness.  PlanarBoundaryCoherence records the
equivalent global-boundary formulation that removes the entire normalization
boundary at once.

\subsection{Research mechanisms and status}

The preceding equivalence theorem is the completed logical dictionary.  The
following mechanisms ask how its divisorial premise might be established
from the special polynomial-plane data; they are not additional steps in the
endgame.

\subsubsection{Two statements and the missing secant--trace morphism}

Put \(T:=A_N\).  For every prime-order subgroup \(C\leq G\), write
\[
  \mathfrak J_C:=\mathfrak a_C+\mathfrak j_C^{\mathrm{mov}},
  \qquad
  \mathcal P_C^{\mathrm{fm}}:=H^1_{\mathfrak J_C}(T).
\]
The ideal \(\mathfrak J_C\) cuts out the scheme-theoretic intersection of
the \(C\)-fixed locus with the boundaries of all sheets moved by \(C\).
Thus \(\mathfrak J_C\subseteq\mathfrak p_E\) is exactly the fixed--moving
height-one configuration isolated by boundary separation.

\begin{proposition}[I. The fixed--moving pole tower]
\label{prop:fixed-moving-pole-tower}
For every height-one prime \(\mathfrak p\subset T\),
\[
  \left(\mathcal P_C^{\mathrm{fm}}\right)_{\mathfrak p}
  \cong
  \begin{cases}
    0,&\mathfrak J_C\nsubseteq\mathfrak p,\\[2pt]
    \Frac(T_{\mathfrak p})/T_{\mathfrak p},
      &\mathfrak J_C\subseteq\mathfrak p.
  \end{cases}
\]
In the second case this localization has principal parts of arbitrarily
large pole order and is not a finite \(T_{\mathfrak p}\)-module.  Moreover,
\[
  \operatorname{Sep}_2(F)
  \quad\Longleftrightarrow\quad
  \mathcal P_C^{\mathrm{fm}}=0
  \quad\text{for every prime-order }C\leq G.
\]
\end{proposition}

\begin{proof}
The action of \(G\) on \(T\) is faithful, so \(\mathfrak J_C\neq0\) for
\(C\neq1\).  Local cohomology commutes with localization.  If
\(\mathfrak J_C\nsubseteq\mathfrak p\), its extension to
\(T_{\mathfrak p}\) is the unit ideal and the localization vanishes.  If
\(\mathfrak J_C\subseteq\mathfrak p\), normality makes
\(T_{\mathfrak p}\) a DVR, and the one-element \v{C}ech complex gives
\[
  H^1_{\mathfrak J_CT_{\mathfrak p}}(T_{\mathfrak p})
  \cong \Frac(T_{\mathfrak p})/T_{\mathfrak p}.
\]
The classes of \(\pi^{-m}\), for a uniformizer \(\pi\) and \(m\geq1\),
have unbounded pole order, proving nonfiniteness.

If boundary separation holds, no height-one prime contains
\(\mathfrak J_C\).  The normal Noetherian surface \(T\) is \(S_2\), so
\(\mathfrak J_C\) has grade at least two and
\(H^1_{\mathfrak J_C}(T)=0\).  Conversely, failure of boundary separation
gives a height-one prime containing \(\mathfrak J_C\), and the displayed
localization shows that \(\mathcal P_C^{\mathrm{fm}}\neq0\).
\end{proof}

\begin{proposition}[II. The trace-dual bounded stage]
\label{prop:trace-dual-bounded-stage}
Define
\[
  T^\dagger
  :=\{z\in N:\operatorname{Tr}_{N/K}(zT)\subseteq B\},
  \qquad
  Q_{\mathrm{tr}}:=T^\dagger/T.
\]
Then the trace pairing gives
\[
  T^\dagger\cong\operatorname{Hom}_B(T,B),
\]
and \(Q_{\mathrm{tr}}\) is a finite \(T\)-module.  If
\(\mathfrak p_E\subset T\) has ramification index \(e_E\) and uniformizer
\(\pi_E\), then
\[
  (Q_{\mathrm{tr}})_{\mathfrak p_E}
  \cong
  \pi_E^{-(e_E-1)}T_{\mathfrak p_E}/T_{\mathfrak p_E}.
\]
Thus the trace-dual quotient contains only one bounded pole stage.
Equivalently, for
\[
  \omega_Z:=\operatorname{Hom}_B(T,\omega_Y),
  \qquad
  \eta_F:=c^{-1}dP\wedge dQ,
\]
finite duality gives
\[
  Q_F^{\omega}
  :=\omega_Z/T\eta_F
  \cong Q_{\mathrm{tr}}.
\]
\end{proposition}

\begin{proof}
The finite locally free \(B\)-algebra \(T\) and generic separability of
\(N/K\) identify \(T^\dagger\) with \(\operatorname{Hom}_B(T,B)\) through
the trace pairing.  Since \(B\) is normal, traces of elements integral over
\(B\) lie in \(B\); hence \(T\subseteq T^\dagger\), and the quotient is
finite.  After height-one localization, characteristic-zero ramification is
tame and the different exponent is \(e_E-1\), giving the displayed
inverse-different formula
\cite{StacksDedekindDifferent,StacksFiniteFlatDifferent}.

Trivializing \(\omega_Y\) by \(dP\wedge dQ\) identifies \(\omega_Z\) with
the trace dual.  The Keller identity identifies \(T\eta_F\) with the
submodule \(T\subseteq T^\dagger\), and therefore identifies the two
quotients \cite{StacksQuasiFiniteDualizing}.
\end{proof}

\paragraph{The trace transporter.}
Fix a prime-order subgroup \(C\leq G\), and put
\[
  U_C:=\Spec T\setminus V(\mathfrak J_C),
  \qquad
  R_C:=\Gamma(U_C,\mathcal O_{U_C})\subseteq N.
\]
Faithfulness of the \(G\)-action gives \(\mathfrak J_C\neq0\), so \(U_C\)
is dense and its ring of regular functions embeds in \(\Frac(T)=N\).
Since \(T\) is a domain, \(H^0_{\mathfrak J_C}(T)=0\), and the affine
open-complement sequence gives
\[
  0\longrightarrow T\longrightarrow R_C
  \longrightarrow H^1_{\mathfrak J_C}(T)\longrightarrow0,
  \qquad
  \mathcal P_C^{\mathrm{fm}}\cong R_C/T
  \cite{StacksLocalCohomologySequence}.
\]
Thus \(T\subseteq T^\dagger\subseteq N\) places both
\(\mathcal P_C^{\mathrm{fm}}\) and \(Q_{\mathrm{tr}}\) in \(N/T\).  Define
the trace-transporter ideal
\[
  \mathfrak c_C^{\mathrm{tr}}
  :=\{s\in T:sR_C\subseteq T^\dagger\}.
\]
For \(s\in T\), the following conditions are equivalent:
\begin{enumerate}[label=\textup{(\roman*)}]
  \item \(s\in\mathfrak c_C^{\mathrm{tr}}\);
  \item
  \[
    \operatorname{Tr}_{N/K}(srt)\in B
    \qquad(r\in R_C,\ t\in T);
  \]
  \item multiplication by \(s\) on \(N/T\) restricts to the unique
  \(T\)-linear morphism
  \[
    \lambda_{C,s}\colon\mathcal P_C^{\mathrm{fm}}\longrightarrow
    Q_{\mathrm{tr}},
    \qquad [r]\longmapsto[sr].
  \]
\end{enumerate}
Indeed, the first two conditions are equivalent by the definition of
\(T^\dagger\).  Under condition \textup{(i)}, one has
\(sR_C\subseteq T^\dagger\), while \(sT\subseteq T\); hence multiplication
by \(s\) descends to the quotient map in \textup{(iii)}.  This is the exact
landing morphism; it is not required to be injective.

\paragraph{The secant--frame denominator candidate.}
The boundary-compatible secant is indexed by an oriented pair of conjugate
sheets.  For \(g\in G\) and \(\sigma\in C\) with
\(\sigma\notin gHg^{-1}\), put
\[
  M_{F;g,\sigma}:=M_F(g(x),\sigma g(x))\in M_2(N),
  \qquad
  A_{F;g,\sigma}:=\operatorname{adj}(M_{F;g,\sigma}).
\]
The two sheet maps agree after applying \(F\), and they are distinct because
\(\sigma gH\ne gH\).  Each \(U_h\) is dense in the integral scheme \(Z\),
so \(U_g\cap U_{\sigma g}\) is nonempty and
\[
  D_{g,\sigma}:=
  \Gamma(U_g\cap U_{\sigma g},\mathcal O_Z)\subseteq N
\]
is a domain.  Regularity on the two respective sheets defines the maps
\[
  A\rightrightarrows D_{g,\sigma},
  \qquad
  a\longmapsto g(a),
  \qquad
  a\longmapsto\sigma g(a).
\]
Proposition~\ref{prop:universal-moved-collision-factorization} therefore
gives the canonical off-diagonal factorization
\[
  C_F^\circ\longrightarrow
  D_{g,\sigma}.
\]
In particular, every entry of \(M_F\), \(\operatorname{adj}(M_F)\), and
\(\Phi_F\), together with the distinguished row \((-d_2,d_1)\), has a
regular image on the overlap.  We do not assert that the Hilbert--Burch
resolution remains exact after this base change or that the row generates
the overlap's canonical module.  This factorization comes from the general
collision construction; it is not yet the trace/conductor landing for an
arbitrary boundary section.
Proposition~\ref{prop:relative-coevaluation-moved-sheet} supplies, for the
same moved sheet, a finite relative-trace reconstruction.  It likewise does
not bound the pole order of its trace coefficients on \(U_C\): the
\(C\)-invariant component is trace-bearing and can still contain arbitrarily
deep poles.  The missing inclusion below is precisely the assertion that the
Keller-specific overlap data controls that surviving component uniformly.
The extended Keller frame of Proposition
\ref{prop:keller-inverse-jacobian-frame} consists of derivations
\(\widetilde\partial_P,\widetilde\partial_Q\) on \(N\).  They commute:
their commutator vanishes on \(K=\C(P,Q)\), and uniqueness of extension
across the finite separable extension \(N/K\) makes it zero on \(N\).  For
\(\alpha=(a,b)\in\mathbf N^2\), set
\(D^\alpha=\widetilde\partial_P^a
\widetilde\partial_Q^b\), and, for \(r\geq0\), let
\[
  \Xi_{g,\sigma,r}:=
  \left\{
    D^\alpha(M_{F;g,\sigma})_{ij},
    D^\alpha(A_{F;g,\sigma})_{ij}
    :1\leq i,j\leq2,\ |\alpha|\leq r
  \right\}\subset N.
\]
Put
\[
  \Lambda_{C,r}^{\mathrm{sf}}
  :=T+\sum_{\substack{g\in G\\C\nsubseteq gHg^{-1}}}
       \ \sum_{\substack{\sigma\in C\\\sigma\ne1}}
       \ \sum_{\xi\in\Xi_{g,\sigma,r}}T\xi\subseteq N,
  \qquad
  \mathfrak d_{C,r}^{\mathrm{sf}}
  :=\{s\in T:s\Lambda_{C,r}^{\mathrm{sf}}\subseteq T\}.
\]
Because \(C\) has prime order, the condition
\(C\nsubseteq gHg^{-1}\) implies
\(C\cap gHg^{-1}=1\); hence every \(\sigma\ne1\) in the displayed sum moves
the sheet.  Indexing over all such \((g,\sigma)\) avoids choices and remains
finite.  The case \(r=1\) is the first-jet candidate obtained from the
pairwise-overlap secant matrix, its adjugate, and the inverse-Jacobian frame.

\begin{proposition}[Nonzero secant--frame denominator]
\label{prop:nonzero-secant-frame-denominator}
For every \(C\) and every \(r\geq0\),
\[
  \mathfrak d_{C,r}^{\mathrm{sf}}\neq0.
\]
Moreover, every nonzero element of this ideal remains nonzero in
\(T_{\mathfrak p}\) for every prime \(\mathfrak p\subset T\).
\end{proposition}

\begin{proof}
The set of generators in \(\Lambda_{C,r}^{\mathrm{sf}}\) is finite, and
each belongs to \(N=\Frac(T)\).  Clearing their finitely many denominators
produces \(0\neq s\in T\) with
\(s\Lambda_{C,r}^{\mathrm{sf}}\subseteq T\).  The localization assertion
follows because \(T\) is a domain.
\end{proof}

\paragraph{The monogenic trace-comparison order.}
The finite separable extension \(N/K\) has a primitive element
\cite{StacksPrimitiveElement}.  Clearing one
\(B\)-denominator gives an integral primitive generator
\(\alpha_{\mathrm{sec}}\in T\) with
\(N=K(\alpha_{\mathrm{sec}})\).  Let
\[
  R_{\mathrm{sec}}:=B[\alpha_{\mathrm{sec}}]\subseteq T,
  \qquad
  f_{\mathrm{sec}}(Z)
  =Z^d+\gamma_{d-1}Z^{d-1}+\cdots+\gamma_0\in B[Z]
\]
be its monic minimal polynomial, where \(d=[N:K]\) and
\(\gamma_d:=1\).  The coefficients lie in \(B\), because they lie in
\(K\), are integral over \(B\), and \(B\) is normal.  Consequently
\[
  R_{\mathrm{sec}}\cong B[Z]/(f_{\mathrm{sec}}),
  \qquad
  R_{\mathrm{sec}}
  =\bigoplus_{j=0}^{d-1}B\alpha_{\mathrm{sec}}^j,
  \qquad
  \Frac(R_{\mathrm{sec}})=N.
\]
Thus \(R_{\mathrm{sec}}/B\) is finite free and a hypersurface, hence a
complete intersection.  Its normalization in \(N\) is the already-defined
ring \(T\).  Put
\[
  J_{\mathrm{sec}}
  :=f_{\mathrm{sec}}'(\alpha_{\mathrm{sec}}),
  \qquad
  \mathfrak f_{\mathrm{sec}}
  :=(R_{\mathrm{sec}}:_N T)
  =\{a\in N:aT\subseteq R_{\mathrm{sec}}\}.
\]
Separability gives \(J_{\mathrm{sec}}\neq0\).  It is the
hypersurface-Jacobian generator of the different
\cite{StacksHypersurfaceDifferent}.  The ideal
\(\mathfrak f_{\mathrm{sec}}\) is, by definition, the conductor of the
finite normalization \(R_{\mathrm{sec}}\subseteq T\).
Here the subscript records the order's role in the secant--trace comparison;
the element \(\alpha_{\mathrm{sec}}\) is not asserted to be generated by
the secant--frame coefficients, and \(J_{\mathrm{sec}}\) is neither
\(\delta_F\) nor the Keller constant.  No such identification is needed for
the reduction below.

\begin{proposition}[Monogenic Tate--conductor reduction]
\label{prop:monogenic-tate-conductor-reduction}
For \(0\leq j<d\), set
\[
  h_j(Z):=\sum_{k=j+1}^{d}\gamma_kZ^{k-1-j}.
\]
Then
\[
  \frac{f_{\mathrm{sec}}(X)-f_{\mathrm{sec}}(Y)}{X-Y}
  =\sum_{j=0}^{d-1}h_j(X)Y^j,
\]
the elements \(h_j(\alpha_{\mathrm{sec}})\) form a triangular
\(B\)-basis of \(R_{\mathrm{sec}}\),
and, for every \(z\in N\),
\[
  \boxed{
  J_{\mathrm{sec}}z
  =\sum_{j=0}^{d-1}
    h_j(\alpha_{\mathrm{sec}})
    \operatorname{Tr}_{N/K}
      (z\alpha_{\mathrm{sec}}^j).}
  \tag{\(\mathrm{Tate}_{\mathrm{sec}}\)}
\]
Moreover,
\[
  R_{\mathrm{sec}}^\dagger
  :=\{z\in N:
      \operatorname{Tr}_{N/K}(zR_{\mathrm{sec}})\subseteq B\}
  =J_{\mathrm{sec}}^{-1}R_{\mathrm{sec}},
\]
and therefore
\[
  \boxed{
  T^\dagger
  =J_{\mathrm{sec}}^{-1}\mathfrak f_{\mathrm{sec}}.}
  \tag{\(\dagger_{\mathrm{sec}}\)}
\]
\end{proposition}

\begin{proof}
Define the \(B\)-linear Tate functional on the power basis by
\[
  \sigma_{\mathrm{sec}}
  \left(\sum_{j=0}^{d-1}b_j\alpha_{\mathrm{sec}}^j\right)
  :=b_{d-1}.
\]
After tensoring the power basis with \(K\), extend
\(\sigma_{\mathrm{sec}}\) uniquely to the corresponding \(K\)-linear
coefficient functional \(N\to K\), and retain the same notation.
The displayed divided difference is the one-variable
Tate--Scheja--Storch coevaluation tensor.  Its reconstruction and trace
identities are
\[
  u=\sum_{j=0}^{d-1}h_j(\alpha_{\mathrm{sec}})
       \sigma_{\mathrm{sec}}(\alpha_{\mathrm{sec}}^ju),
  \qquad
  \operatorname{Tr}_{N/K}(u)
  =\sigma_{\mathrm{sec}}(J_{\mathrm{sec}}u).
\]
Applying the first identity to \(J_{\mathrm{sec}}z\) and the trace identity
to each \(\alpha_{\mathrm{sec}}^jz\) gives
\((\mathrm{Tate}_{\mathrm{sec}})\).  Equivalently,
\(J_{\mathrm{sec}}^{-1}h_j(\alpha_{\mathrm{sec}})\) is trace-dual to
\(\alpha_{\mathrm{sec}}^j\).  This proves
\(R_{\mathrm{sec}}^\dagger
=J_{\mathrm{sec}}^{-1}R_{\mathrm{sec}}\)
\cite{SchejaStorch1975,StacksTateMap,StacksDedekindDifferent}.

Finally,
\[
  T^\dagger
  =(R_{\mathrm{sec}}^\dagger:_N T)
  =(J_{\mathrm{sec}}^{-1}R_{\mathrm{sec}}:_N T)
  =J_{\mathrm{sec}}^{-1}(R_{\mathrm{sec}}:_N T),
\]
which is \((\dagger_{\mathrm{sec}})\).
Indeed, the first equality follows directly from
\[
  zT\subseteq R_{\mathrm{sec}}^\dagger
  \quad\Longleftrightarrow\quad
  \operatorname{Tr}_{N/K}(zTR_{\mathrm{sec}})\subseteq B
  \quad\Longleftrightarrow\quad
  \operatorname{Tr}_{N/K}(zT)\subseteq B,
\]
using \(R_{\mathrm{sec}}\subseteq T\) and
\(1\in R_{\mathrm{sec}}\).
\end{proof}

\paragraph{Missing morphism: Secant--Frame Trace Landing.}
The explicit first-jet research target is the ideal containment
\[
  \boxed{
  \ \mathfrak d_{C,1}^{\mathrm{sf}}
  \subseteq\mathfrak c_C^{\mathrm{tr}}\ }
  \tag{\(\mathrm{SFTL}_C\)}
\]
for every prime-order \(C\leq G\).  Equivalently, every
secant--frame denominator must satisfy the uniform trace identity
\[
  \operatorname{Tr}_{N/K}(srt)\in B
  \qquad(s\in\mathfrak d_{C,1}^{\mathrm{sf}},\ r\in R_C,\ t\in T).
\]
Proposition~\ref{prop:monogenic-tate-conductor-reduction} gives the
equivalent conductor form
\[
  \boxed{
  J_{\mathrm{sec}}\,
  \mathfrak d_{C,1}^{\mathrm{sf}}R_C
  \subseteq\mathfrak f_{\mathrm{sec}}}
  \quad\Longleftrightarrow\quad
  \mathfrak d_{C,1}^{\mathrm{sf}}
  \subseteq\mathfrak c_C^{\mathrm{tr}}.
  \tag{\(\mathrm{SCTL}_C\)}
\]
Because the left-hand side is already a \(T\)-submodule, containment in
the conductor is equivalently
\[
  J_{\mathrm{sec}}\,
  \mathfrak d_{C,1}^{\mathrm{sf}}R_C
  \subseteq R_{\mathrm{sec}}.
\]

This also makes the finite part of the reduction explicit.  Since \(T\) is
finite over \(B\), it is finite over \(R_{\mathrm{sec}}\).  Choose
\(R_{\mathrm{sec}}\)-module generators
\(t_1,\ldots,t_m\) of \(T\).  For \(s\in T\) and \(r\in R_C\),
\[
\begin{aligned}
  sr\in T^\dagger
  &\Longleftrightarrow
  J_{\mathrm{sec}}sr\in\mathfrak f_{\mathrm{sec}}\\
  &\Longleftrightarrow
  J_{\mathrm{sec}}srt_i\in R_{\mathrm{sec}}
    \quad(1\leq i\leq m)\\
  &\Longleftrightarrow
  \operatorname{Tr}_{N/K}
    (srt_i\alpha_{\mathrm{sec}}^j)\in B
    \quad(1\leq i\leq m,\ 0\leq j<d).
\end{aligned}
\]
The final equivalence follows from
\((\mathrm{Tate}_{\mathrm{sec}})\) and trace duality.  Thus Tate reduction
replaces the universal test element \(t\in T\) by one fixed finite family of
trace coordinates.  It does not prove those coordinates integral uniformly
for the arbitrary boundary section \(r\in R_C\); that uniform-in-\(r\)
conductor landing is precisely the remaining morphism.

The relative coevaluation isolates the exact part still uncontrolled.  If
\(p=|C|\) and
\[
  \operatorname{Av}_C(z):=\frac1p\sum_{\sigma\in C}\sigma(z)\in N^C,
\]
then transitivity of trace gives
\[
  \operatorname{Tr}_{N/K}(z)
  =\operatorname{Tr}_{N^C/K}
      \bigl(p\,\operatorname{Av}_C(z)\bigr).
\]
Thus the finite family above is equivalently the family of assertions
\[
  \operatorname{Tr}_{N^C/K}\!\left(
    p\,\operatorname{Av}_C
      (srt_i\alpha_{\mathrm{sec}}^j)\right)\in B.
\]
The non-invariant character components do not contribute to the trace.  The
remaining problem is therefore a uniform integrality theorem for this
\(C\)-invariant remainder.  This is substantive: in the standard tame local
model \(\sigma(\pi)=\zeta\pi\), the invariant poles
\(\pi^{-p},\pi^{-2p},\ldots\) survive averaging.  Hence neither cyclic
averaging nor the finite relative reconstruction alone supplies the missing
pole bound.

\begin{proposition}[Scalar landing collapses the boundary overring]
\label{prop:scalar-landing-collapses-overring}
If \(0\neq s\in\mathfrak c_C^{\mathrm{tr}}\), then
\[
  R_C=T,
  \qquad
  \mathcal P_C^{\mathrm{fm}}=0.
\]
Equivalently, the same conclusion follows from a nonzero \(s\in T\) such
that
\[
  J_{\mathrm{sec}}sR_C\subseteq R_{\mathrm{sec}}.
\]
In particular no height-one prime of \(T\) contains
\(\mathfrak J_C\).
\end{proposition}

\begin{proof}
The conductor identity converts \(sR_C\subseteq T^\dagger\) into
\(J_{\mathrm{sec}}sR_C\subseteq\mathfrak f_{\mathrm{sec}}\), hence into
the displayed containment in \(R_{\mathrm{sec}}\subseteq T\).  Put
\(a=J_{\mathrm{sec}}s\neq0\).  For every \(r\in R_C\) and every
\(n\geq0\), the ring property of \(R_C\) gives \(r^n\in R_C\), and
therefore
\[
  ar^n\in T.
\]
Thus every \(r\in R_C\subseteq\Frac(T)\) is almost integral over \(T\).
A Noetherian domain has the property that almost integral elements of its
fraction field are integral; normality then places them in the domain
\cite{StacksAlmostIntegral}.  Hence \(R_C\subseteq T\).  The reverse
inclusion is built into the open-complement sequence, so \(R_C=T\) and
\(\mathcal P_C^{\mathrm{fm}}=R_C/T=0\).  Statement~I now gives the final
height-one assertion.
\end{proof}

Assuming \((\mathrm{SFTL}_C)\), the nonzero-denominator proposition supplies
\(0\neq s_C^{\mathrm{sf}}\in\mathfrak d_{C,1}^{\mathrm{sf}}\subseteq
\mathfrak c_C^{\mathrm{tr}}\).  Multiplication then defines the required map
explicitly by
\[
  \lambda_{C,s_C^{\mathrm{sf}}}([r])=[s_C^{\mathrm{sf}}r].
\]
Proposition~\ref{prop:scalar-landing-collapses-overring} then gives the
global chain
\[
  0\neq s_C^{\mathrm{sf}},\quad
  J_{\mathrm{sec}}s_C^{\mathrm{sf}}R_C\subseteq R_{\mathrm{sec}}
  \quad\Longrightarrow\quad
  R_C=T
  \quad\Longrightarrow\quad
  \mathcal P_C^{\mathrm{fm}}=0.
\]
The local bounded-pole calculation below remains the geometric explanation
for why this is impossible at a hidden divisor, but is not needed for this
global implication.

Merely asking for a nonzero element of
\(\mathfrak d_{C,1}^{\mathrm{sf}}\cap\mathfrak c_C^{\mathrm{tr}}\) would
not define a new mechanism: since \(T\) is a domain and
\(\mathfrak d_{C,1}^{\mathrm{sf}}\neq0\), that intersection is nonzero
exactly when \(\mathfrak c_C^{\mathrm{tr}}\neq0\), which is boundary
separation itself.  In fact the dichotomy below shows that the full
containment \((\mathrm{SFTL}_C)\) is logically equivalent to the same
separation assertion.  Its value as a research target is therefore not a
weaker logical premise, but the prescribed coefficientwise trace identity
that one can try to prove directly from the displayed secant--frame data.
No identity proved here yet shows that first jets suffice; jet order (1) is
the first explicitly typed candidate.

More generally, write
\((\mathrm{SFTL}_C)=(\mathrm{SFTL}_{C,1})\), and for every finite
\(\ell\geq0\) state
\[
  (\mathrm{SFTL}_{C,\ell})\colon
  \mathfrak d_{C,\ell}^{\mathrm{sf}}
  \subseteq\mathfrak c_C^{\mathrm{tr}}.
\]
Any one finite-order statement suffices by the same argument, whereas
\((\mathrm{SFTL}_{C,1})\) is the designated first-jet proposal.  A minimal scalar
certificate would instead specify, by an explicit secant--frame formula, one
\(0\neq s_{C,\ell}^{\mathrm{sf}}\in
\mathfrak d_{C,\ell}^{\mathrm{sf}}\) and prove its trace landing.  Merely
asserting that some nonzero element lies in the intersection is not a new
mechanism; the formula producing that element is essential.  The ideal form
used above is choice-free and asks for landing coefficientwise.

The nonvanishing proposition is only finite denominator clearing; it does
not prove landing for the generally nonfinite overring \(R_C\).  Indeed, if
a height-one prime \(\mathfrak p\) contains \(\mathfrak J_C\), Statement I
and the localized open-complement sequence, together with Statement II, give
\[
  (R_C)_{\mathfrak p}=\Frac(T_{\mathfrak p}),
  \qquad
  (T^\dagger)_{\mathfrak p}
  =\pi^{-(e_{\mathfrak p}-1)}T_{\mathfrak p}.
\]
For every \(0\neq s\in T\), multiplication by \(s/1\) on
\(\Frac(T_{\mathfrak p})/T_{\mathfrak p}\) is surjective and therefore
cannot land in the proper bounded stage
\(\pi^{-(e_{\mathfrak p}-1)}T_{\mathfrak p}/T_{\mathfrak p}\).
Consequently \((\mathfrak c_C^{\mathrm{tr}})_{\mathfrak p}=0\).  Since
\(T\to T_{\mathfrak p}\) is injective, this forces
\(\mathfrak c_C^{\mathrm{tr}}=0\), and
\((\mathrm{SFTL}_C)\) excludes \(\mathfrak p\).  Notice that no unit or
avoidance condition \(s_C^{\mathrm{sf}}\notin\mathfrak p\) is needed:
global nonzeroness already survives localization in the domain \(T\).

Conversely, boundary separation gives \(R_C=T\), hence
\(\mathfrak c_C^{\mathrm{tr}}=T\).  Thus nonzero trace landing is an exact
Keller-specific formulation of the open divisorial assertion, rather than
a consequence of finite flatness.  The secant and frame data now provide a
concrete finite candidate ideal; the missing theorem is the uniform passage
from those finitely many coefficients to every \(r\in R_C\).

Thus the purity argument remains the unique endgame.  Statements I and II
identify its unbounded source and finite target, while
\((\mathrm{SFTL}_C)\) is a concrete first-jet construction problem whose
solution would establish its height-one premise.

\begin{leanfiles}
The stable modules
\leanpath{CollisionIdeals/Planar/Boundary/Separation.lean} and
\leanpath{CollisionIdeals/Planar/Rigidity/Consequences.lean} formalize the
canonical ideals \leanpath{fixedLocusIdeal},
\leanpath{movingBoundaryIdeal}, and \leanpath{fixedMovingBoundaryIdeal},
together with \leanpath{PlanarBoundarySeparation}, the pulled-back
conjugate boundaries, their valuation-theoretic containment at ramified
divisors, and the hypothesis-parametrized purity endgame.  That endgame
takes branch purity, finite-\'etale rigidity, and a supplied
\leanpath{PlanarKellerCollisionModel F} explicitly.  The proposition types
are in \leanpath{CollisionIdeals/Planar/External/Interfaces.lean};
\leanpath{CollisionIdeals/Planar/External/Assumptions.lean} supplies default
branch-purity, finite-\'etale, and planar Ax--Grothendieck witnesses as
axioms.  Ax--Grothendieck is needed only for the terminal automorphism
statement, not for obstruction vanishing.  The Lean separation predicate
currently quantifies over every nontrivial subgroup; formalizing its
prime-order Cauchy reduction is a remaining routine API task.

The boundary object \leanpath{BoundaryPrincipalParts} is defined in
\leanpath{CollisionIdeals/Planar/Boundary/PrincipalParts.lean}; its canonical
specialization \leanpath{FixedMovingBoundaryPrincipalParts} is in
\leanpath{CollisionIdeals/Planar/Research/FixedMovingBoundaryPrincipalParts.lean}.
The local-cohomology localization, DVR computation, and \(S_2\)-vanishing
used in Statement~I are not yet formalized.  The research module
\leanpath{CollisionIdeals/Planar/Research/PrincipalPartsStrategy.lean}
defines \leanpath{TraceIntegralSubmodule} and
\leanpath{boundedStageComparisonIdeal}; it proves the former equal to
Mathlib's \leanpath{Submodule.traceDual} on the unit order.  The opt-in
modules \leanpath{CollisionIdeals/Planar/Research/TateReconstruction.lean}
and \leanpath{CollisionIdeals/Planar/Research/MonogenicTraceDual.lean}
formalize the finite Tate sum, the monogenic trace-dual formula, and the
overorder conductor--codifferent identity.  Trace-dual/Hom finiteness and
the local inverse-different assertion of Statement~II remain unfinished in
the planar specialization.

Secant evaluation on a generic collision map pair and its canonical
factorization through the off-diagonal quotient are formalized in
\leanpath{CollisionIdeals/Planar/ConjugateSecantEvaluation.lean}.  The actual
Galois \((g,\sigma g)\) map pair and its planar off-diagonal specialization
are in
\leanpath{CollisionIdeals/General/Galois/PolynomialCollisionPair.lean} and
\leanpath{CollisionIdeals/Planar/GaloisCollisionPair.lean}; the relative-
\(C\) coevaluation, overlap-section-ring specialization, and evaluated
coefficient family remain missing.  The choice-free telescoping
coefficients and polynomial Keller frame are in
\leanpath{CollisionIdeals/Planar/ExplicitSecant.lean} and
\leanpath{CollisionIdeals/Planar/KellerFrame.lean}.  The theorem
\leanpath{finiteCoefficientDenominatorIdeal_ne_bot} in
\leanpath{CollisionIdeals/Planar/Research/SecantFrameDenominator.lean}
formalizes finite denominator clearing, but the actual conjugate evaluated
family is missing.

The opt-in module
\leanpath{CollisionIdeals/Planar/Research/MonogenicOrder.lean} directly
reuses Mathlib's \leanpath{Algebra.adjoin},
\leanpath{Algebra.adjoin.powerBasis'}, \leanpath{minpoly.equivAdjoin}, and
\leanpath{conductor} for the monogenic comparison order.  The module
\leanpath{CollisionIdeals/Planar/Research/MonogenicLanding.lean} defines the
conductor bounded stage and its transporter and proves an abstract
landing-to-endgame implication from a supplied pole witness.  The
open-section algebra \(R_C\), its common embedding with the bounded stage in
\(N/T\), the concrete conjugate evaluated secant--frame family, the
specialization of the generic finite-overring theorem to \(R_C\), and the
ideal containment \((\mathrm{SFTL}_C)\) remain to be formalized.  These
fixed-moving, trace, denominator, and monogenic landing modules remain
behind \leanpath{CollisionIdeals/Planar/Research.lean}; the secant, frame,
Galois-pair, conjugate-evaluation, and basic boundary modules named above
belong to the stable \leanpath{CollisionIdeals/Planar.lean} import spine.
\end{leanfiles}

\begin{remark}[Exact remaining content of classical \(JC(2)\)]
\label{rem:exact-remaining-jc2}
The traditional complex conjecture \(JC(2)\) has no boundary, coverage, or
module-finiteness condition in its definition: it asserts simply that every
planar Keller map is a polynomial automorphism.  Corollary
\ref{cor:planar-jc2-boundary-forms} proves that its universal content is
equivalent to the boundary formulations above.  Theorem
\ref{thm:planar-galois-rigidity} settles every normal marked extension, and
Corollary~\ref{cor:quadratic-planar-rigidity} settles generic degree at most
two.  Consequently, to prove \(JC(2)\) it remains to show that every planar
Keller map with \(L/K\) nonnormal and \(d_F\geq3\) satisfies
\[
  \begin{gathered}
    \operatorname{codim}_Z
    V\!\left(\mathfrak a_C+\mathfrak j_C^{\mathrm{mov}}\right)\geq2
    \quad\text{for every prime-order }C\leq\Gal(N/K).
  \end{gathered}
\]
Equivalently, one must exclude a common height-one fixed--moving boundary
component in the normal closure of every remaining nonnormal extension.
This is an exact
reformulation of the unresolved part of \(JC(2)\), not an additional
hypothesis silently built into the classical conjecture.
\end{remark}

\section{\texorpdfstring{The cubic \(S_3\) obstruction in dimension three}
  {The cubic S3 obstruction in dimension three}}
\label{sec:cubic-s3-obstruction}

We now specialize the dimension-independent constructions of
Section~\ref{sec:general-construction} to dimension three and isolate the
separable nonnormal cubic case.

Let \(S_3=\operatorname{Perm}(\{1,2,3\})\) be the symmetric group on
three letters.

\begin{proposition}[Dimension-three counterexample criterion]
\label{prop:dimension-three-counterexample-criterion}
For a polynomial map
\[
  F\colon\A^3_{\C}\longrightarrow\A^3_{\C},
\]
consider the following statements:
\begin{enumerate}[label=\textup{(\arabic*)}]
  \item \(\det JF\in\C^\times\);
  \item \(F\) is not a polynomial automorphism;
  \item \(R_F^\circ\neq\varnothing\).
\end{enumerate}
Statements \textup{(2)} and \textup{(3)} are equivalent.  Consequently,
\(F\) is a dimension-three Jacobian-conjecture counterexample exactly
when \textup{(1)} and \textup{(3)} hold.
\end{proposition}

\begin{proof}
By Corollary~\ref{cor:off-diagonal-empty-criterion} and
Theorem~\ref{thm:automorphism-criterion},
\[
  R_F^\circ=\varnothing
  \quad\Longleftrightarrow\quad
  I_R=I_\Delta
  \quad\Longleftrightarrow\quad
  F\in\Aut_{\mathrm{poly}}(\A^3_{\C}).
\]
Negating this equivalence proves
\textup{(2)}\(\Longleftrightarrow\)\textup{(3)}.  Combining this with
\textup{(1)} gives the final assertion.
\end{proof}

\subsection{The cubic extension and its normal closure}
\label{sec:cubic-extension}

Let
\[
  F=(F_1,F_2,F_3)\colon\A^3_{\C}\longrightarrow\A^3_{\C}
\]
be dominant and generically finite.  Let
\[
  A=\C[x_1,x_2,x_3],
  \qquad
  B=\C[F_1,F_2,F_3]\subseteq A,
  \qquad
  K=\Frac(B),
  \qquad
  L=\Frac(A),
\]
and let
\[
  \theta\colon K\otimes_BA\longrightarrow L,
  \qquad
  \lambda\otimes a\longmapsto\lambda a,
\]
be the isomorphism of
Theorem~\ref{thm:generic-collision-base-change}.  Suppose that \(L/K\)
is separable and nonnormal, and
choose a power basis with generator \(\alpha\) and degree three:
\[
  L=K(\alpha),
  \qquad
  [L:K]=3.
\]
If \(m_\alpha\in K[T]\) is the minimal polynomial of \(\alpha\), let
\(h_\alpha\in L[T]\) be determined by
\[
  m_\alpha(T)=(T-\alpha)h_\alpha(T)
  \qquad\text{in }L[T].
\]
Let \(N/K\) be the normal closure of \(L/K\), containing \(L\) through
the marked inclusion.  Let
\[
  G=\Gal(N/K),
  \qquad
  H=\Gal(N/L).
\]
Recall from Section~\ref{sec:normalized-covers-inertia} that
\[
  G/H\xrightarrow{\sim}\Hom_K(L,N),
  \qquad
  gH\longmapsto g|_L,
  \qquad
  \core_G(H)=1.
\]
The final equality follows because \(N\) is the normal closure of
\(L/K\); equivalently, the action of \(G\) on the conjugate embeddings
of \(L\) is faithful.

For \(Y=\Spec B\), the finite-normalization construction of
Section~\ref{sec:normalized-covers-inertia} gives
\[
  Z=\Norm_N(Y)\longrightarrow
  \overline X=\Norm_L(Y)\longrightarrow Y.
\]
If \(E\in Z^{(1)}\), with decomposition group \(D_E\) and inertia group
\(I_E\), then its conjugate centers are indexed by
\(\mathcal Q_E=D_E\backslash G/H\).  For
\(\omega=D_EgH\), Proposition~\ref{prop:intermediate-ramification-index}
states
\[
  e\bigl(\operatorname{cent}(E,\omega)/\nu(E)\bigr)
  =
  i_E(\omega)
  =
  [I_E:I_E\cap gHg^{-1}].
\]

The following proposition is the standard residual-factor description
of a separable nonnormal cubic extension.  We include its short proof
to fix the marked identification with \(N\).

\begin{proposition}[Cubic residual-factor identification]
\label{prop:cubic-tensor-decomposition}
The polynomial \(h_\alpha\) is irreducible over \(L\), and there is an
\(L\)-algebra isomorphism
\[
  \varphi\colon L[T]/(h_\alpha)\xrightarrow{\sim}_L N.
\]
Consequently,
\[
  \Psi\colon L\otimes_KL\xrightarrow{\sim}_L L\times N,
  \qquad
  \operatorname{pr}_1\circ\Psi=\mu_L.
\]
\end{proposition}

\begin{proof}
Write
\[
  m_\alpha(T)=T^3+a_2T^2+a_1T+a_0.
\]
If the quadratic \(h_\alpha\) were reducible over \(L\), it would have a
root \(\beta\in L\).  Separability gives \(\beta\neq\alpha\), and the
third root would be
\[
  \gamma=-a_2-\alpha-\beta\in L.
\]
Thus \(m_\alpha\) would split over \(L\).  Since
\(L=K(\alpha)\), the extension \(L/K\) would then be the splitting
field of the separable polynomial \(m_\alpha\), hence normal.  This
contradicts the assumed nonnormality, so \(h_\alpha\) is irreducible.

Let \(\beta\) denote the image of \(T\) in \(L[T]/(h_\alpha)\).  This
field contains the three roots
\[
  \alpha,\qquad \beta,\qquad -a_2-\alpha-\beta
\]
of \(m_\alpha\), and is generated over \(L\) by \(\beta\).  It is
therefore the splitting field of \(m_\alpha\), which is the normal
closure \(N\) of \(L/K\).  Choose the resulting \(L\)-algebra
isomorphism
\[
  \varphi\colon L[T]/(h_\alpha)\xrightarrow{\sim}_L N.
\]

Let
\[
  \Phi_\alpha\colon
  L\otimes_KL
  \xrightarrow{\sim}_L
  L\times L[T]/(h_\alpha)
\]
be the marked-root decomposition of
Theorem~\ref{thm:marked-root-tensor-decomposition}.  It satisfies
\[
  \operatorname{pr}_1\circ\Phi_\alpha=\mu_L.
\]
Composing its second factor with \(\varphi\), let
\[
  \Psi=(\id_L\times\varphi)\circ\Phi_\alpha.
\]
Since \(\id_L\times\varphi\) leaves the first factor unchanged,
\(\operatorname{pr}_1\circ\Psi=\mu_L\).
\end{proof}

\begin{theorem}[Galois group of a nonnormal cubic extension]
\label{thm:index-three-core-free-s3}
\[
  \Gal(N/K)\cong S_3.
\]
\end{theorem}

\begin{proof}
Proposition~\ref{prop:cubic-tensor-decomposition} identifies \(N\) with
the degree-two extension \(L[T]/(h_\alpha)\) of \(L\).  Hence
\[
  |H|=[N:L]=2.
\]
The normal-closure identities recalled above give
\[
  \core_G(H)=1,
  \qquad
  [G:H]=[L:K]=3.
\]
Hence the action of \(G\) on the three cosets of \(H\) gives
\[
  \rho\colon G\longrightarrow\operatorname{Perm}(G/H)\cong S_3,
  \qquad
  \ker\rho=\core_G(H)=1.
\]
Core-freeness therefore makes \(\rho\) injective, so
\[
  |G|\leq |S_3|=6.
\]
The index formula gives
\[
  |G|=[G:H]|H|=3\cdot2=6.
\]
Thus \(|G|=6\), and the injective homomorphism \(\rho\) is an
isomorphism \(G\cong S_3\).
\end{proof}

\subsection{The cubic off-diagonal factor}

Retain the cubic data of Section~\ref{sec:cubic-extension}, and suppose
in addition that \(F\) is Keller.  Thus
\[
  B=\C[F_1,F_2,F_3]\subseteq A=\C[x_1,x_2,x_3],
  \qquad
  K=\Frac(B)\subseteq L=\Frac(A)=K(\alpha)\subseteq N,
\]
where \(L/K\) is separable, nonnormal, and of degree three, and \(N/K\)
is its normal closure.  Let
\[
  \theta\colon K\otimes_BA\xrightarrow{\sim}L,
  \qquad
  \lambda\otimes a\longmapsto\lambda a.
\]

\begin{corollary}[Generic cubic \(S_3\) collision]
\label{cor:cubic-s3-collision}
There is a \(K\)-algebra isomorphism
\[
  \Psi_F\colon
  K\otimes_BC_F\xrightarrow{\sim}_K L\times N,
  \qquad
  \operatorname{pr}_1\circ\Psi_F
  =
  \theta\circ(1\otimes\bar\mu_F).
\]
It satisfies
\[
  \Psi_F\!\left(
    \ker\bigl(\theta\circ(1\otimes\bar\mu_F)\bigr)
  \right)
  =
  0\times N,
  \qquad
  \Gal(N/K)\cong S_3.
\]
Consequently,
\[
  \Obs(F)\neq0,
  \qquad
  I_R\subsetneq I_\Delta,
  \qquad
  R_F^\circ\neq\varnothing,
  \qquad
  F\notin\Aut_{\mathrm{poly}}(\A^3_{\C}).
\]
\end{corollary}

\begin{proof}
Theorem~\ref{thm:generic-collision-base-change} gives a \(K\)-algebra
isomorphism
\[
  \Phi_F\colon K\otimes_BC_F\xrightarrow{\sim}_K L\otimes_KL
\]
such that
\[
  \mu_L\circ\Phi_F
  =
  \theta\circ(1\otimes\bar\mu_F).
\]
Proposition~\ref{prop:cubic-tensor-decomposition} gives
\[
  \Psi\colon L\otimes_KL\xrightarrow{\sim}_L L\times N,
  \qquad
  \operatorname{pr}_1\circ\Psi=\mu_L.
\]
After restricting \(\Psi\) from \(L\) to \(K\), let
\[
  \Psi_F=\Psi\circ\Phi_F.
\]
It follows that
\[
  \operatorname{pr}_1\circ\Psi_F
  =
  \theta\circ(1\otimes\bar\mu_F),
  \qquad
  \Psi_F\!\left(
    \ker\bigl(\theta\circ(1\otimes\bar\mu_F)\bigr)
  \right)
  =
  0\times N\neq0.
\]

Suppose that the diagonal-evaluation homomorphism
\(\bar\mu_F\colon C_F\to A\) were injective.  Since \(K\) is a
localization of \(B\), it is flat over
\(B\), so \(1\otimes\bar\mu_F\) would also be injective.
Theorem~\ref{thm:generic-collision-base-change} shows that \(\theta\)
is an isomorphism.  Therefore
\(\theta\circ(1\otimes\bar\mu_F)\) would be injective, contradicting
the nonzero kernel \(0\times N\) above.  Hence
\(\ker(\bar\mu_F)\neq0\), and
Proposition~\ref{prop:obstruction-kernel} gives
\(\Obs(F)\neq0\).  Proposition
\ref{prop:collision-contained-in-diagonal} shows that
\(I_R\subseteq I_\Delta\) strict, and
Corollary~\ref{cor:off-diagonal-empty-criterion} gives
\(R_F^\circ\neq\varnothing\).

Finally, Theorem~\ref{thm:index-three-core-free-s3} identifies
\(\Gal(N/K)\) with \(S_3\).  If \(F\) were a polynomial automorphism,
Theorem~\ref{thm:automorphism-criterion} would give
\(\Obs(F)=0\), contrary to \(\Obs(F)\neq0\).  Since \(F\) is Keller, it is
a dimension-three Jacobian-conjecture counterexample.
\end{proof}

For the same cubic map and field tower, let
\[
  X=\Spec A,
  \qquad
  Y=\Spec B,
  \qquad
  G=\Gal(N/K),
  \qquad
  H=\Gal(N/L).
\]
Suppose there is a diagram over \(Y\)
\[
  \mathcal D=
  \left(
    Z=\Norm_N(Y)\xrightarrow{\pi}
    \overline X=\Norm_L(Y)\xrightarrow{\bar\nu}Y,\,
    j\colon X\hookrightarrow\overline X
  \right)
\]
in which \(\pi\) and \(\bar\nu\) are finite, \(j\) is an open immersion,
\(\bar\nu\circ j=f_F\), and \(f_F\colon X\to Y\) is \'etale.  Let
\[
  \partial X=\overline X\setminus j(X),
  \qquad
  G\cong S_3
\]
by Theorem~\ref{thm:index-three-core-free-s3}.

\begin{corollary}[Hidden inertia in the \(S_3\)-normalization]
\label{cor:cubic-hidden-inertia}
\[
  \mathcal H(\mathcal D)=\Ram^{(1)}(Z/Y).
\]
\end{corollary}

\begin{proof}
Apply Proposition~\ref{prop:etale-sheet-hidden-inertia} to the displayed
normalization diagram.
\end{proof}

\begin{remark}[Ordered roots]
Let
\[
  m_\alpha(T)=T^3+a_2T^2+a_1T+a_0
\]
have roots \(\alpha_1,\alpha_2,\alpha_3\) in \(N\).  For distinct
\(i,j,k\),
\[
  \alpha_k=-a_2-\alpha_i-\alpha_j,
  \qquad
  K(\alpha_i,\alpha_j)
  =
  K(\alpha_1,\alpha_2,\alpha_3)
  =
  N.
\]
Thus the off-diagonal factor, which marks two distinct roots, is already
the ordered-root normal closure.
\end{remark}

\begin{leanfiles}
\leanpath{CollisionIdeals/ComplexThree/Cubic/FunctionField.lean},
\leanpath{CollisionIdeals/ComplexThree/Cubic/GaloisGroup.lean},
\leanpath{CollisionIdeals/ComplexThree/Cubic/S3Collision.lean},
\leanpath{CollisionIdeals/ComplexThree/Cubic/Branch.lean},
\leanpath{CollisionIdeals/ComplexThree/OffDiagonal.lean},
\leanpath{CollisionIdeals/ComplexThree/Cubic/Main.lean}, and
\leanpath{CollisionIdeals/ComplexThree/Statements/JacobianConjecture.lean}.  The
lower-level generic Lean witness builder takes the residual-field
isomorphism and \(H\neq1\) as explicit inputs.  The cubic branch module now
derives both from nonnormality, and the dimension-three wrapper uses those
derived inputs.  Hidden inertia is formalized separately, pointwise at a
supplied ramified divisor and conditional on a supplied Keller
collision--normalization model and \(H\neq1\), in
\leanpath{CollisionIdeals/ComplexThree/Cubic/HiddenInertia.lean}.
\end{leanfiles}

\begin{remark}[Exact relation with classical complex \(JC(3)\)]
\label{rem:exact-relation-jc3}
The traditional complex three-dimensional Jacobian conjecture asserts that
every polynomial map
\[
  F\colon\A^3_{\C}\longrightarrow\A^3_{\C},
  \qquad
  \det JF\in\C^\times,
\]
is a polynomial automorphism; equivalently, every such map satisfies
\(\Obs(F)=0\).  Corollary~\ref{cor:cubic-s3-collision} proves the following
implication under a prescribed function-field hypothesis:
\[
  \begin{gathered}
    \det JF\in\C^\times,\qquad [L:K]=3,\qquad
    L/K\text{ nonnormal}\\
    \Longrightarrow\quad \Obs(F)\neq0
    \quad\Longrightarrow\quad
    F\text{ is a counterexample to }JC(3).
  \end{gathered}
\]
Consequently, the existence of one such Keller map would disprove
\(JC(3)\), while \(JC(3)\) implies that no such map exists.  This section
does not construct such a map, and it does not assert that every possible
counterexample to \(JC(3)\) has generic degree three.  It identifies the
generic collision geometry of this cubic nonnormal subclass.
\end{remark}

\section{Conclusion}
\label{sec:conclusion}

The obstruction ideal \(\Obs(F)\) vanishes exactly when the
scheme-theoretic collision relation \(X\times_YX\) equals the diagonal.
In dimension two there is one endgame with several exact formulations.
PlanarBoundaryCoherence expresses it through disappearance of the whole
deleted boundary.  PlanarRamificationRigidity, moving-sheet coverage, and
boundary separation express its height-one premise in differential,
valuation-theoretic, and finite-group language.  Only after that premise is
established do purity
and finite-\'etale rigidity make the normal closure trivial and force
\(\Obs(F)=0\).  The image of the deleted boundary is exactly the
nonproperness set, and the boundary divisor-class argument proves the
normal/Galois case unconditionally.  Universal validity of the equivalent
planar criteria is exactly \(JC(2)\), not a result asserted here.  In
the cubic dimension-three class, by contrast, the nonzero generic factor is
the ordered-root \(S_3\)-normal closure.
Both statements arise from the same collision, off-diagonal, and
conjugate-sheet constructions; their point is to locate the obstruction,
respectively, at the normalization boundary and on the generic
off-diagonal sheet.

Under PlanarBoundaryCoherence or PlanarRamificationRigidity in dimension
two, under the conditions of Section~\ref{sec:cubic-s3-obstruction} in
dimension three, and with
\(\Ram^{(1)}(Z/Y)\neq\varnothing\) in the cubic normalization, the
contrast is
\[
  \begin{array}{c@{\qquad}|@{\qquad}c}
    \hline
    \text{\rm planar rigidity criteria}
      & \text{\rm nonnormal cubic class in dimension three}\\ \hline
    \operatorname{PBC}(F)\text{\rm\ or }\operatorname{PRR}(F)
      & \text{\rm hidden inertia in the \(S_3\)-cover realized}\\
    \Obs(F)=0 & \Obs(F)\neq0\\
    I_R=I_\Delta & I_R\subsetneq I_\Delta\\
    F\text{\rm\ is a polynomial automorphism}
      & F\text{\rm\ is not a polynomial automorphism}.\\
    \hline
  \end{array}
\]

\paragraph{Motivating example.}
The explicit three-dimensional polynomial counterexample announced by
Levent Alpöge \cite{Alpoge2026} and independently formalized in
Isabelle/HOL by Ramos, Hulak, and de Queiroz
\cite{AFPJacobianCounterexample} motivated the cubic \(S_3\) comparison.
The formal verification establishes a
constant nonzero Jacobian determinant and distinct rational points with
the same image.  That particular map is not used in the arguments above:
the dimension-three result is the structural implication from the
stated cubic and marked-normal-closure conditions.

The formalization uses Lean~4 and Mathlib \cite{Lean4,Mathlib}.

\paragraph{Acknowledgment.}
OpenAI Codex (GPT-5.6 Sol) assisted with Lean proof engineering,
mathematical and bibliographic checking, and manuscript editing.  The
author reviewed and takes responsibility for all statements, proofs,
citations, and formalization code.

\bibliographystyle{alpha}
\bibliography{references}

\end{document}
